\documentclass[12pt]{cours}

\newtheorem{openpb}{Open Problem}

\title{Algorithmics}
\author{Pierre Aboulker, Paul Jeanmaire, Tatiana Starikovskaya}
\date{\today}

\begin{document}
\part[First Algorithms]{Lecture 1 - 28/09}
\localtableofcontents
\section{Organisation}
Mail Tatiana : \url{starikovskaya@di.ens.fr}
Homeworks are 30\% of the final grade, final (theory from lecture)
Textbooks :
\begin{itemize}
    \item \textsl{Introduction to Algorithms} - Cormen, Leiserson, Riverst, Stein
    \item \textsl{Algorithms on strings, trees, and sequences} - Gusfield
    \item \textsl{Approximation Algorithms} - Vazirani
    \item \textsl{Parametrized Algorithms} - Cygan, Fomin, Kowalik, Lokshtanov, Marx, Pilipczuk, Saurabh
\end{itemize}

\section{Introduction}
Algorithm take Inputs and give an output.
\begin{openpb}[Mersenne Prime]
    Find a new prime of form $2^{n} - 1$
\end{openpb}
Algorithms do not depend on the language. Algorithms should be simple, fast to write and efficient.
Word RAM model : Two Parts : one with a constant number of registers of $w$ bits with direct access, and one with any number of registers, only with indirect access (pointers).
Allows for elementary operations: basic arithmetic and bitwise operations on registers, conditionals, goto, copying registers, halt and malloc.
To index the memory storing input of size $n$ with $n$ words, we need register length to verify $w \geq \log n$
Algorithms can always be rewritten using only elementary operations.
Complexity :
\begin{itemize}
    \item $Space(n)$ is the maximum number of memory words used for input of size $n$
    \item $Time(n)$ is the maximum number of \textsl{elementary} operations used for input of size $n$
\end{itemize}
Complexity Notations :
\begin{itemize}
    \item $f \in \mathcal{O}(g)$ if $\exists n_0 \in \mathbb{N}, c \in \mathbb{R}_{+},\ f(n) \leq c \cdot g(n), \ \forall n \geq n_0$
    \item $f \in \Omega(g)$ if $\exists n_0 \in \mathbb{N}, c \in \mathbb{R}_{+},\ f(n) \geq c \cdot g(n), \ \forall n \geq n_0$
    \item $f \in \Theta(g)$ if $\exists n_0 \in \mathbb{N}, c_{1}, c_{2} \in \mathbb{R}_{+},\ c_{1} \cdot g(n) \leq f(n) \leq c_{2} \cdot g(n), \ \forall n \geq n_0$
\end{itemize}

\section{Data Structures}
\subsection{Introduction}
Way to store elements of a data base that is created to answer frequently asked queries using pre-processing.
We care about space used, construction, query and update time.
Can be viewed as an algorithm, which analysed on basics.
Containers are basic Data Structures, maintaining the following operations :
\begin{enumerate}
    \item Random Access : given $i$, access $e_{i}$
    \item Access first/last element
    \item Insert an element anywher
    \item Delete any element
\end{enumerate}

\subsection{Array}
An array is a pre-allocated contiguous memory area of a \emph{fixed} size. It has random access in $\mathcal{O}(1)$, but doesn't allow insertion nor deletion.

Linear Search : given an integer $x$ return $1$ if $e_{i} = x$ else 0.
\begin{algorithm}
    \caption{Linear Search in an Array. \\ Complexity : Time = $\mathcal{O}(n)$ | Space = $\mathcal{O}(n)$}
\end{algorithm}

\subsection{Doubly Linked List}
Memory area that does not have to be contiguous and consists of registers containing a value and two pointers to the previous and next elements.
It has random access in $\mathcal{O}(n)$, access/insertion/deletion at head/tail in $\mathcal{O}(1)$.
\begin{algorithm}
    \caption{Insertion in a Doubly Linked List \\ Complexity : $\mathcal{O}(1)$}
    \begin{algorithmic}
        \Input
        $ L, x$
        \EndInput
        \State $x.next \gets L.head$
        \If{$L.head \neq NIL$}
        \State $L.head.prev \gets x$
        \EndIf
        \State $L.head \gets x$
        \State $x.prev = Nil$
    \end{algorithmic}
\end{algorithm}


\subsection{Stack and Queue}
Stack : Last-In-First-Out data structure, abstract data structure. Access/insertion/deletion to top in $\mathcal{O}(1)$.

\begin{openpb}[Optimum Stack Generation]
    Given a finite alphabet $\Sigma$ and $X \in \Sigma^{n}$. Find a shortest sequence of stack operations push, pop, emit that prints out $X$. You must start and finish with an empty stack.
    Current best solution is in $\tilde{\mathcal{O}}(n^{2.8603})$.
\end{openpb}

Queue : First-In-First-Out abstract data structure. Access to front, back in $\mathcal{O}(1)$, deletion and insertion at front and back in $\mathcal{O}(1)$.

\section{Approaches to algorithm design}
Solve small sub-problems to solve a large one.

\subsection{Dynamic Programming}
Break the problem into many closely related sub-problems, memorize the result of the sub-problems to avoid repeated computation

Examples :
\begin{algorithm}
    \caption{Recursive Fibonacci Numbers \\ Complexity: Exponential}
    \begin{algorithmic}
        \State RFibo($n$) :
        \Input
        $ n$
        \EndInput
        \If {$n \leq 1$}

        \Return $n$
        \EndIf\\
        \Return RFibo($n-1$) + RFibo($n-2$)
    \end{algorithmic}
\end{algorithm}

\begin{algorithm}
    \caption{Dynamic Programming Fibonacci Numbers \\ Time = $\mathcal{O}(n)$ | Space = $\mathcal{O}(n)$}
    \begin{algorithmic}
        \Input
        $ n$
        \EndInput   
        \State $Tab \gets zeros(n)$ \Comment $zeros(n)$ returns a $n$-array of zeros.
        \State $Tab[0] \gets 0$
        \State $Tab[1] \gets 1$
        \For{$i \gets 2$ to $n$}
        \State $Tab[i] = Tab[i-1] + Tab[i-2]$
        \EndFor\\
        \Return Tab[n]
    \end{algorithmic}
\end{algorithm}

Levenshtein Distance between two strings can be computed in $\mathcal{O}(mn)$ instead of exponential time. Based on \url{https://arxiv.org/pdf/1412.0348.pdf}, this is the best one can do.
RNA folding : retrieving the 3D shape of RNA based on their representation as strings. Currently, we know it is possible to find $\mathcal{O}(n^3)$, in $\tilde{\mathcal{O}}(n^{2.8606})$ and if \emph{SETH} is true, there is no $\mathcal{O}(n^{\omega - \epsilon})$. We know $\omega \in \left[2, 2.3703\right)$

\begin{openpb}
    Is there a better Complexity for RNA folding ? What is the true value of $\omega$ ?
\end{openpb}

Knapsack problem : An optimization problem with bruteforce complexity $\mathcal{O}(2^{n})$.
\begin{algorithm}
    \caption{Knapsack : Dynamic Programming \\
        Time = $\mathcal{O}(nW)$ | Space = $\mathcal{O}(nW)$}
    \begin{algorithmic}
        \Input $\ W, w, v$ \Comment Capacity, weight and values vectors.
        \EndInput

        \State $KP = zeros(n, W)$

        \For {$i \gets 0$ to $n$}
        \State $KP[i, 0] = 0$
        \EndFor

        \For {$w \gets 0$ to $W$}
        \State $KP[0, w] = 0$
        \EndFor

        \For {$i \gets 0$ to $n$}
        \For {$w \gets 0$ to $W$}
        \If {$w < w_{i}$}
        \State $KP[i, w] \gets KP[i-1, w]$
        \Else
        \State $KP[i, w] = \max\begin{cases}
                 & KP[i-1, w]                 \\
                 & KP[i-1, w - w_{i}] + v_{i}
            \end{cases}$.
        \EndIf
        \EndFor
        \EndFor\\
        \Return $KP[n, W]$
    \end{algorithmic}
\end{algorithm}

\subsection{Greedy Techniques}
Problems solvable with the greedy technique form a subset of those solvable with DP. Problems must have the optimal substrcture property.
Principle : choosing the best at the moment.


Example : The Fractional Knapsack Problem\\
Algorithm : Iteratively select the greatest value-per-weight ratio. \\
\begin{theorem}
    This algorithm returns the best solution, in time $\mathcal{O}(n \log n)$
\end{theorem}
\begin{proof}[By contradiction]
    Suppose we have $\frac{v_{1}}{w_{1}} \geq \ldots \geq \frac{v_{n}}{w_{n}}$. Let $ALG = p = (p_{1}, \ldots, p_{n})$ be the output by the algorithm and $OPT = q = (q_{1}, \ldots, q_{n})$ be optimal.\\
    Assume that $OPT \neq ALG$, let $i$ be the smallesst index such $p_{i} \neq q_{i}$. There is $p_{i} > q_{i}$ by construct. \\
    Thus, there exists $j > i$ such that $p_{j} < q_{j}$. We set $q' = (q_{1}^{'}, \ldots, q_{n}^{'}) = (q_{1}, \ldots, q_{i-1}, q_{i} + \epsilon, q_{i+1}, \ldots, q_{j} - \epsilon\frac{w_{i}}{w_{j}}, \ldots, q_{n})$\\
    $q^{'}$ is a feasible solution : $\sum\limits_{i = 1}^{n} q_{i}^{'} \cdot w_{i} = \sum\limits_{i = 1}^{n} q_{i} w_{i} \leq W$\\
    Yet, $\sum\limits_{i = 1}^{n} q_{i}^{'} \cdot v_{i} > \sum\limits_{i = 1}^{n} q_{i} \cdot v_{i}$, ce qui contredit la
\end{proof}

\part[Recursion]{Lecture 2 - 5/10}
\localtableofcontents
\section{Divide and Conquer}
Divide a problem into smaller ones to solve those, then combine the solutions to get a solution of the bigger problem.\\
Example : \textsl{Merge Sort} : Its complexity verifies $T(n) = T(\lceil n/2\rceil) + T(\lfloor n/2 \rfloor) + \O(n)$. From that we will derive that $T(n) = \O(n\log n)$

\section{Analysis of Recursive Algorithms}
We have recurrences we want to solve. We have three methods :


\subsection{Substitution Method}
The method :
\begin{enumerate}
    \item Guess the asymptotic of $T(n)$
    \item Show the answer via induction
\end{enumerate}
For \textsl{Merge Sort}: we guess $T(n) \leq c \cdot n\log_{2}n, \forall n \geq 2$. We choose $c$ that verifies that property until $n = 6$. \\
Substituting in the recurrence equation : \[T(n) \leq cn\log_{2}\frac{n}{2} + c\log_{2}\frac{n}{2} + c\frac{n+2}{2} + a\cdot n = c \dot n\log_{2}n + a \cdot n + c \cdot \log_{2}n - c\frac{n}{2}\]
If we then choose $c$ so that it is bigger than $20a$ we get :
\[T(n) \leq cn\log_{2}n + a\cdot n - c\cdot n/20 \leq cn\log_{2}n\]

\subsection{Recursion-tree Method}
\begin{enumerate}
    \item Simplify the equation :
          \begin{itemize}
              \item Delete floors and ceils
              \item Suppose $n$ is of a good form
          \end{itemize}
    \item Draw a tree, rooted with the added term and the recursive calls
    \item Start again, and recursively fill the tree
\end{enumerate}
We get a tree of depth $\log_{k}n$ if $n$ is divided by $k$ in recursive calls. We can now sum the values of the nodes, to get an approximation, and start verifying.

\section{Master Theorem}
\subsection{The Theorem}
\begin{theorem}[Master Theorem]\label{thm:MT}
    If we have recurrence equation $T(n) = aT(n/b) + f(n)$ where $a \geq 1, b > 1$ are integers, $f(n)$ is asymptotically positive. Let $r = \log_{b}a$, we have :
    \begin{enumerate}
        \item If $f(n) = \O(n^{r-\epsilon})$ for some $\epsilon > 0$, then $T(n) = \Theta(n^r)$
        \item If $f(n) = \Theta(n^{r})$ then $T(n) = \Theta(n^{r}\log n)$
        \item If $f(n) = \Omega(n^{r + \epsilon})$ for some $\epsilon > 0$, and $af(n/b) \leq cf(n)$ for some constant $c < 1$ and all sufficiently large $n$ (regularity condition) then $T(n) = \Theta(f(n))$.
    \end{enumerate}
\end{theorem}
\begin{remark}
    The Master Theorem \ref{thm:MT} does not cover all possible cases for $f(n)$. Example : $f(h) = h^{r}/\log h$
\end{remark}
\begin{remark}
    The Master Theorem \ref{thm:MT} is sometimes called \textsc{Théorème sur les Récurrences de Partition}
\end{remark}

\section{The Proof}
Plan : \begin{itemize}
    \item Analyse the recurrence as if $T$ is defined over reals (continuous version)
    \item Prove the discrete version
\end{itemize}

\subsection{Continuous Master Theorem}
\begin{proof}
    \begin{lemma}
        Define $T(n) = \begin{cases}
                \Theta(1)      & \text{if } n \leq \hat{n} \\
                aT(n/b) + f(n) & \text{if } n > \hat{n}
            \end{cases}$
        Then \[T(n) = \Theta\left(n^{r}\right) + \sum_{k=0}^{\lceil \log_{b}(n/\hat{n})\rceil - 1}a^{k}f(n/b^{k})\]
    \end{lemma}
    \begin{proof}
        In the Recursion-Tree, stopped when the argument of $T$ is smaller than $\hat{n}$ which is when depth is $\lceil \log_{b}(n/\hat{n})\rceil - 1$, we get :
        \[
            \begin{aligned}
                T(n) & \leq \sum_{k=0}^{\lceil \log_{b}(n/\hat{n})\rceil - 1}a^{k}f(n/b^{k}) + \Theta(a^{log_{b}(n/\hat{n})}) \\ &= \sum_{k=0}^{\lceil \log_{b}(n/\hat{n})\rceil - 1}a^{k}f(n/b^{k}) + \Theta(a^{\log_{b}(n)})\\ &= \sum_{k=0}^{\lceil \log_{b}(n/\hat{n})\rceil - 1}a^{k}f(n/b^{k}) + \Theta(n^{\log_{b}(a)})
            \end{aligned}
        \]
    \end{proof}
    Back to the proof :
    \begin{lemma}
        Define $g(n) = \Theta(n^{r}) + \sum_{k= 0}^{q} a^{k}f(n/b^{k})$ Then :
        \begin{enumerate}
            \item If $f(n) = \O(n^{r - \epsilon})$ then $g(n) = \Theta(n^{r})$
            \item If $f(n) = \Theta(n^{r})$ then $g(n) = \Theta(n^{r}\log{n})$
            \item If $f(n) = \Omega(n^{r + \epsilon})$ and we have the regularity condition then $g = \Theta(f)$
        \end{enumerate}
    \end{lemma}
    \begin{proof}
        \begin{enumerate}
            \item Case 1 : \[
                      g(n)
                      \begin{aligned}
                           & = \Theta(n^{r}) + \sum_{k= 0}^{q} a^{k}f(n/b^{k})                              \\
                           & = \Theta(n^{r}) + \O\left(\sum_{k= 0}^{q} a^{k}(n/b^{k})^{r - \epsilon}\right) \\
                      \end{aligned}
                  \]
                  However :
                  \[
                      \begin{aligned}
                          \sum_{k= 0}^{q} a^{k}(n/b^{k})^{r - \epsilon} & = n^{r-\epsilon}\sum_{k = 0}^{q} (ab^{\epsilon}/b^{r})^{k}                                                                      \\
                                                                        & = n^{r-\epsilon}\sum_{k=0}^{\lceil \log_{b}(n/\hat{n})\rceil - 1}(b^{\epsilon})^{k} = \O(n^{r - \epsilon}(n/\hat{n})^{epsilon})
                      \end{aligned}
                  \]
                  Thus : $g(n) = \Theta(n^{r})$
                  % --------------------------------------------------------------------------
            \item Case 2 : We have :
                  \[
                      \begin{aligned}
                          g(n) & = \Theta(n^{r}) + \sum_{k=0}^{q} a^{k}f(n/b^{k})                      \\
                               & = \Theta(n^{r}) + \Theta\left(\sum_{k=0}^{q}a^{k}(n/b^{k})^{r}\right) \\
                      \end{aligned}
                  \]
                  However :
                  \[
                      \begin{aligned}
                          \sum_{k=0}^{q}a^{k}(n/b^{k})^{r} & = n^{r}\sum_{k = 0}^{q} (a/b^{r})^{k}                                                        \\
                                                           & = n^{r}\sum_{k = 0}^{\lceil \log_{b}(n/\hat{n})\rceil - 1}1 = n^{r}\Theta(\log_{b}n/\hat{n})
                      \end{aligned}
                  \]
                  % ----------------------------------------------------------------------------
            \item Case 3 : By induction on $k$ : $a^{k}f(n/b^{k}) \leq c^{k}f(n)$. Thus :
                  \[
                      \sum_{k = 0}^{q}a^{k}f(n/b^{k}) \leq \sum_{k = 0}^{q}c^{k}f(n) = f(n) \sum_{k= 0}^{q}c^{k} = \Theta(f(n))
                  \]
        \end{enumerate}
    \end{proof}
    We thus have proved the continuous Master Theorem.
\end{proof}

\subsection{Discrete Master Theorem}
We have now showed the continuous Master Theorem, following \textsc{William Kuszmaul, Charles E. Leiserson}, \href{https://epubs.siam.org/doi/pdf/10.1137/1.9781611976496.15}{\textit{Floors and Ceilings in Divide-and-Conquer Recurrences}}, Symposium on Simplicity in Algorithms 2021.\\
\begin{proof}
    See slides below
    \includepdf[nup = 1x2, pages = 32-45]{Lecture_2.pdf}
\end{proof}
\begin{remark}[Remarks on the Proof]
    \begin{itemize}
        \item Lemmas 1 to 3 serve to show that the argument does not go too far when it is rounded up or down.
        \item Slide 36 Last Line : $\frac{2}{\beta^{i} - 1} < \frac{4}{\beta^{i}}$ for $i \geq i_{0}$. Thus : $\sum_{i = 1}^{\infty} \frac{2}{\beta^{i} - 1} < \sum_{i = 1}^{\infty}\frac{4}{\beta^{i}} + \sum_{i = 0}^{i_0}\frac{2}{\beta^{i} - 1}$ and that last sum is $\O(1)$
        \item Slide 37 Line 3 : The first inequalities comes from the Recursion-Tree, so that we can ensure the argument does not deviate to much, by the second inequalities.
    \end{itemize}
\end{remark}
\subsection{Use Cases}
Using the Master Theorem we can show the complexity of many algorithms :
\begin{enumerate}
    \item Merge Sort Complexity : $T(n) = T(\lceil n/2 \rceil) + T(\lfloor n/2 \rfloor) + \O(n) = \Theta(n\log n)$
    \item Strassen's Algorithm for Matrix Multiplication : $T(n) = 7 T(n/2) + \Theta(n^{2}) \Rightarrow T(n) = \O(n^{\log_{2}7}) = \O(n^{2.8074})$
\end{enumerate}

\section{Fast Multiplication of Polynomials}
The sum of two degree $n$ polynomials can be done in $\O(n)$, Horner's rule for evaluation produces $\O(n)$ complexity. The naïve product can be done in $\O(n^2)$\\
Remembering Lagrange's Theorem on Polynomials (or Vandermonde's Determinant, or anything really), degree $n$ polynomials are entirely represented by their point-value reprensentation over $n$ distinct points $(x_{i}, y_{i})$.
Then, by converting the coefficient reprensentation to point-value representation, then by point-wise multiplicating the polynomials, then by going back to the coefficient representation, we can have a better algorithm.

\subsection{Point-Value Multiplication}
It is easily done in $\O(n)$ if both polynomials are represented over the same axis.

\subsection{Coefficient to Point-Value Conversion - Fast Fourier Transform}
For $P = \sum_{i = 0}^{n-1} a_{i}x^{i}$, we define : \[
    \begin{aligned}
         & P_{odd}(x)  & = a_{n-1}x^{n/2-1} + a_{n-3}x^{n/2-3} + \ldots + a_{1}x              \\
         & P_{even}(x) & = a_{n-2}x^{n/2-2} + a_{n-4}x^{n/2-4} + \ldots+ a_{2}x^{2/2} + a_{0}
    \end{aligned}
\]
\begin{enumerate}
    \item We have : $P = xP_{odd}(x^{2}) + P_{even}(x^{2})$
    \item We evaluate $P_{odd}, P_{even}$ at $(\omega_{n}^{i})^{2}$ recursively by the halving property.
    \item We combine the result.
\end{enumerate}

\subsection{Point-Value to Coefficient Conversion - Inverse Fast Fourier Transform}
\begin{theorem}
    $V_{n}^{-1}[i, j] = \omega_{n}^{-ij}/n$
\end{theorem}

\part[Hashing]{Lecture 3 - 12/1}
\localtableofcontents

Answering queries of appartenance on a set, maintaining a dictionary. Python dictionaries do not have upper bound guarantees, we shouldn't use them. \\
Suppose we have a set $S$ over a universe $U$. We dnote by $n$ the number of keys, $m$ the size of the hash table

\section{Naïve Array-based implementation}
Here we assume that no two objects have equal keys. We store $1$ in a bitvector of length $|U|$ at each position $i$ such that $i$ is in $S$.\\
This takes $\O(|U|)$ space, for $\O(1)$ search time, and $\O(1)$ modification time.

\section{Chained Hash Tables}
We give ourselves a function $h : U \rightarrow \left[1, m\right]$. We send the keys to their image by $h$ in a table. However, since we have no guarantee that $\left|U\right| \leq m$, there might be collisions.\\
To deal with that, instead of storing in the table a boolean, we store a list of all the keys corresponding to $h(k)$. \\
Then we insert a key in constant time $\O(1)$, search a key in time $\O(\left|h[key]\right|)$ and delete a key in time $\O(\left|h[key]\right|)$ since in the worst case the key is at the end of the list.\\

Further analysis leads to the following theorem.
\begin{theorem}[Simple Uniform Hashing Assumption]
    Assuming SUHA : "\textit{$h$ equally distributes the keys into the table slots}", and assuming $h(x)$ can be computed in $\O(1)$, $E\left[T_{search}(n)\right] = \O\left(1 + \frac{n}{m}\right)$, and same for deletion time. \\
    Formally, SUHA is :
    \[
        \begin{aligned}
            \forall y \in \left[1, \abs{T}\right]                & \mathbb{P}\left(h(x) = y\right) = \frac{1}{\abs{T}}                                \\
            \forall y_{1}, y_{2} \in \left[1, \abs{T}\right]^{2} & \mathbb{P}\left(h(x_{1}) = y_{1},\ h(x_{2}) = y_{2}\right) = \frac{1}{\abs{t}^{2}}
        \end{aligned}
    \]
\end{theorem}
\begin{proof}
    \begin{enumerate}
        \item Unsuccessful Search : Suppose that $k_{0}, \ldots, k_{n-1}$ are keys in the dictionary and we perform an unsuccessful search for a key $k$.
              The number of comparisons is : $\sum_{i = 0}^{n-1} \mathbb{1}_{h(k) = h(k_{i})}$. Then, the expected time is : $\mathbb{E}[T_{search}] = \frac{n}{\abs{T}}$ by SUHA.
        \item Successful Search : Suppose keys were introduced in order $k_{0}, \ldots, k_{n-1}$. $k_{i}$ appears before any of $k_{0}, \ldots, k_{i-1}$ and after any of $k_{i+1}, \ldots, k_{n-1}$ that are in the same linked list. Then, to search for $k_{i}$, we need $\sum_{j = i + 1}^{n- 1} \mathbb{1}_{h(k_{j}) = h(k_{i})}$. Under SUHA, the expectation of each of these variables is $\frac{1}{\abs{T}}$. Then, the average expected search time is : $\frac{1}{n}\sum_{i = 0}^{n-1} = 1 + \frac{1}{n}\sum_{i = 0}^{n-1}\frac{(n-1-i)}{\abs{T}} = \O(1 + \frac{n}{\abs{T}})$
    \end{enumerate}
    This concludes the proof of the theorem.
\end{proof}

Good hash functions are functions that distribute the keys evenly. Yet, we do not know what the keys are, and thus will need various heuristics to answer this question. At least, without loss of generality, we can assume that keys are integers.

\subsection{Heuristic Hash Functions}
\begin{itemize}
    \item Division Method : $h(k) = k \mod m$. It is better to choose $m$ to be a prime number, and avoid $m = 2^{p}$.
    \item Multiplication Method : $h(k) = \left\lfloor m \left\{k\cdot A\right\}\right\rfloor$. $A \in (0, 1)$ and $m = 2^{p}$.
\end{itemize}
Yet, fixing the function can allow anyone to construct a probability distribution for which the function will be "bad".

\subsection{Universal Family of Hash Functions}
$H = \left\{h : U \rightarrow \left[0, \abs{T}-1\right]\right\}$ is Universal if :
\[
    \forall k_{1} \neq k_{2} \in U, \ \abs{\left\{h \in H \mid h(k_{1}) = h(k_{2})\right\}} \leq \frac{\abs{H}}{m}
\]
\begin{theorem}
    If $h$ is a hash function chosen uniformly at random from a universal family of hash functions. Suppose that $h(k)$ can be computed in constant time and there are at most $n$ keys. Then the expected search time is $\O(1 + \frac{n}{\abs{T}})$
\end{theorem}
\begin{proof}
    Same as the case when $h$ satisfies SUHA. Observe that the probability comes this time from choosing the function.
\end{proof}

\begin{theorem}
    Let $p \in \mathcal{P}$ such that $U \subseteq \left[0, p-1\right]$. Then \[H = \left\{h_{a, b}(k) = \left(\left(ak + b\right) \mod p\right) \mod \abs{T} \mid a \in \Z_{p}^{*} b \in \Z_{p} \right\}\] is a universal family.
\end{theorem}
\begin{proof}
    Let $k_{1} \neq k_{2} \in U$. Let $l_{i} = (ak_{i} + b) \mod p$. We have $l_{1} \neq l_{2}$ and $ a = ((l_{1} - l_{2})((k_{1}-k_{2})^{-1} \mod p) \mod p) \text{ and } b = (l_{1} - ak_{1}) \mod p$. There is then one-to-one mapping between $(a, b)$ and $(l_{1}, l_{2})$. The number of $h \in H$ such that $h(k_{1}) = h(k_{2})$ is :
    \[
        \abs{\left\{(l_{1}, l_{2})\mid l_{1} \neq l_{2} \in \Z_{p}, \ l_{1} = l_{2} \mod m \right\}} \leq \frac{p(p-1)}{\abs{T}} \leq \frac{\abs{H}}{\abs{T}}
    \]
\end{proof}


\section{Open Addressing}
Elements are stored in the table. To insert, we probe the hash table until we find $x$ or an empty slot. If we find an empty slot, insert $x$ here. To define which slots to probe, we use a hash function that depends on the key and the probe number. To search, we probe the hash table until we either find $x$ (return YES) or an empty slot (return NO). In the analysis, we will assume $h$ to be uniform.

\begin{theorem}[Analysis]
    Given an open-address hash-table with load factor $\alpha = \frac{n}{\abs{T}} < 1$, the expected number of probes in an unsuccessfulsearch is at most $\frac{1}{1-\alpha}$, assuming uniform hashing.
\end{theorem}
\begin{proof}
    An unsuccessful search on $x$ means that every probed slot except the last one is occupied and does not contain $x$, and the last one is empty.
    We define $A_{i}$ the event " The $i$-th probe occurs and is occupied.". By Bayes's Theorem, we must estimate :
    \[
        \mathbb{P}[\text{\# of probes} \geq i] = \prod_{k = 1}^{i - 1} \mathbb{P}[A_{k} \mid \bigcap\limits_{j = 1}^{k - 1} A_{j}]
    \]
    But we have : $\mathbb{P}[A_{j}\mid A_{1} \cap \ldots \cap A_{j - 1}] = \frac{n - j + 2}{\abs{T} - j + 2}$
    So we have : $\mathbb{P}[\text{\# of probes} \geq i] \leq \frac{n}{\abs{T}}^{i - 1} = \alpha^{i - 1}$
    Then, the expected number of probes is :
    \[
        \sum_{i = 1}^{+\infty} \mathbb{P}[\text{\# of probes} \geq i] \leq \sum_{i = 1}^{+\infty} \alpha^{i - 1} = \frac{1}{1 - \alpha}
    \]
\end{proof}
\begin{corollary}
    The expected number of probes during insertion is at most $\frac{1}{1 - \alpha}$.
\end{corollary}
\begin{proof}
    If we insert $x$ we first ran an unsuccessful search for it
\end{proof}

\begin{theorem}
    The expected number of probes during a Successful search is at most $\frac{1}{\alpha}\ln\frac{1}{1- \alpha}$.
\end{theorem}
\begin{proof}
    A successful search for $x$ probes the same sequence of slots as insertion.\\
    If $x$ is the $i$-th element inserted into the table, insertion probes less than $\frac{1}{1 - \frac{i}{\abs{T}}}$ slots in expectation.
    Therefore, the expected time of a successful search is at most :
    \[
        \frac{1}{n}\sum_{i = 0}^{n - 1}\frac{\abs{T}}{\abs{T} - i} = \frac{\abs{T}}{n}\sum_{i = 0}^{n - 1}\frac{1}{\abs{T} - i} = \sum_{i = 0}^{\abs{T}}\frac{1}{i} - \sum_{i = 0}^{\abs{T}-n} \frac{1}{i} \leq \frac{\abs{T}}{n}\ln\frac{\abs{T}}{\abs{T} - n} = \frac{1}{\alpha}\ln\frac{1}{1 - \alpha}
    \]
\end{proof}

This is hard to implement, so we will use heuristics.

\subsection{Heuristic hash functions}
Let $h^{'}, h^{''}$ be two auxiliary hash functions.
\begin{itemize}
    \item Linear Probing : $h(k, i) = (h^{'}(k) + i) \mod \abs{T}$
          This is easy to implement but it suffers from clustering.
    \item Quadratic Probing : $h(k, i) = (h^{'}(k) + c_{1}i + c_{2}i^{2}) \mod \abs{T}$
          We must choose the constants $c_{1}$ and $c_{2}$ carefully, and this still suffers from clustering.
    \item Double Hashing : $h(k, i) = (h^{'}(k) + i h^{''}(k)) \mod \abs{T}$
          To use the whole table, $h^{''}(k)$ must be relatively prime to $m$, e.g. $h^{''}(k)$ is always odd, $m = 2^{i}$.
\end{itemize}

\section{Cuckoo Hashing}
Hashing scheme with search time constant in the worst case, as it maintains a hash function without collisions to achieve perfect hashing. This is possible if the set of keys is static.
Assume that we have two hash functions $h_{1}, h_{2}$ that satisfy SUHA. We store $x$ in either $T[h_{1}(x)]$ or $T[h_{2}(x)]$. Search for $x$ is done in constant time.

\begin{algorithmic}
    \Function{Insert}x
    \If {$ x = T[h_{1}(x)] \text{ or } x = T[h_{2}(x)]$}
    \Return
    \EndIf

    \State $pos \gets h_{1}(x)$
    \For {$i \gets 0 \text{ to } n$}
    \If $T[pos] = Null$
    \State $T[pos] = x$; \Return
    \EndIf

    \State $x \longleftrightarrow T[pos]$
    \If {$pos = h_{1}(x)$}
    \State $pos \gets h_{2}(x)$
    \Else
    \State $pos \gets h_{1}(x)$
    \EndIf
    \EndFor
    \State \textsc{Rehash}
    \State \textsc{Insert}(x)
    \EndFunction
\end{algorithmic}
\begin{theorem}
    This insertion is done in constant time.
\end{theorem}
\begin{lemma}
    Suppose that $\abs{T} \geq c \cdot n$ for some $c > 1$. For any $i, j$, the probability that there exists a path from $i$ to $j$ of length $l \geq 1$ which is a shortest path from $i$ to $j$ is at most $\frac{1}{c^{l}\cdot \abs{T}}$
\end{lemma}
\begin{proof}
    By induction on $l$ :
    \begin{itemize}
        \item Initialization : By SUHA, $\mathbb{P}[h_{1, 2}(x) = y] = \frac{1}{\abs{T}}$. Thus $\mathbb{P}[\text{there is an edge from }i \text{ to } j] = \frac{n}{\abs{T}^{2}} \leq \frac{1}{c\abs{T}}$
        \item Heredity : For $l \geq 1$ there must exist $k$ such that there is a path of length $l - 1$ from $i$ to $k$ and an edge from $k$ to $j$.
    \end{itemize}
\end{proof}
\begin{proof}
    We define the bucket of $x$ as all the cells that can be reached either from $h_{1}(x)$ or $h_{2}(x)$. $x, y$ are in the same bucket if and only if there is a path from $\left\{h_{1}(x), h_{2}(x)\right\}$ to $\left\{h_{1}(y), h_{2}(y)\right\}$. Then :
    \[
        \mathbb{P}[x, y \text{ are in the same bucket}] \leq 4 \sum_{l = 1}^{\infty}\frac{1}{c^{l}\abs{T}} = \frac{4}{(c - 1)\abs{T}}
    \]
    So :
    \[
        \mathbb{E}[\abs{\text{bucket of } x}] = \mathbb{E}\sum X_{x, y} = \sum \mathbb{E}[X_{x, y}] = \sum \mathbb{P}[x, y \text{ are in the same bucket}] \leq \frac{4}{c - 1}
    \]
    Hence, in the absence of rehash, expected insertion time in constant.\\

    The probability that we need a rehash is at most probability that there is a cycle, i.e. a path from $i$ to $i$ : $\frac{4}{c - 1} \leq \frac{1}{2}$.
    The probability that we will need $d$ rehashes is then at most $\frac{1}{2^{d}}$.
    Thus, the expected time per insertion is :
    \[
        \frac{1}{n}\cdot\O(n)\sum_{d = 1}^{+\infty}\frac{1}{2^{d}} = \O(1)
    \]
\end{proof}

\subsection{Rolling Hash Functions}
\begin{definition}
    The Karp-Rabin fingerprint of a string $S = s_{1}\ldots s_{m}$ is :
    \[
        \phi(s_{1}\ldots s_{m}) = \sum_{i = 1}^{m}s_{i}r^{m - 1}\mod p
    \]
    where $p$ is a big enough prime and $r \in \mathbb{F}_{p}$.
\end{definition}
\begin{proposition}
    \begin{itemize}
        \item If S = T, then $\phi(S) = \phi(T)$
        \item Else, $\phi(S) \neq \phi(T)$ with high probability.
    \end{itemize}
\end{proposition}
\begin{proof}
    Let $\sigma$ be the size of the alphabet, $p \geq max \left\{\sigma, n^{c}\right\}$ where $c > 1$ is a constant. We have:
    \[
        \phi(S) = \phi(T) \Leftrightarrow \sum_{i = 1}^{m}(s_{i} - t_{i}) \cdot r^{m - i} \mod p = 0
    \]
    Hence, $r$ is a root of $P(x) = \sum_{i = 1}^{m}(s_{i} - t_{i}) \cdot x^{m-i}$ a polynomial over $\mathbb{F}_{p}$. The probability of such event is at most $\frac{m}{p} \leq \frac{1}{n^{c-1}}$.
\end{proof}
The Karp-Rabin algorithm is as follows :
\begin{itemize}
    \item Compute the fingerprint of the pattern
    \item Compare it with the fingerprint of each $m$-length substring of the text. If the fingerprint is equal to the fingerprint of a substring, report it as an occurrence
\end{itemize}
\begin{proposition}
    We have :
    \[
        \phi(s_{1}\ldots s_{j+1}) = \phi(s_{1}\ldots s_{j})\cdot r + s_{j+1} \mod p
    \]
\end{proposition}
\begin{proof}
    Observe that :
    \[
        \begin{aligned}
            \phi(s_{1}\ldots s_{j+1}) & = \sum_{i = 1}^{j+1}s_{i}r^{j+1-i} \mod p \\
            \
        \end{aligned}
    \]
\end{proof}

\part[Integer Sets]{Lecture 4 : 26/10}
\localtableofcontents
\section{Binary Search Trees}
\begin{definition}
    A Binary Tree is a rooted tree with every node having at most two children. Binary search trees are binary trees verifying the following :
    \begin{center}
        If a node contains element $l$, the left subtree rooted at the left child contains only elements $\leq l$ and the right subtree elements $> l$.
    \end{center}
\end{definition}

\begin{proposition}
    To access a given element, the time needed is $\O(h)$ time where $h$ is the height of the tree.
\end{proposition}
$h$ is not necessarily small, there are many implementations to bound $h$.

\subsection{Red-Black Trees}
\begin{definition}
    A red-black tree is a binary tree that satisfies the following red-black properties :
    \begin{enumerate}
        \item Every node has a color : red or black
        \item The root is black
        \item Every leaf (NIL) is black
        \item If a node is red, both its children are black
        \item For each node, all paths from the node to descendant leaves contain the same number of black nodes.
    \end{enumerate}
\end{definition}
\begin{lemma}
    The height of a red-black tree with $n$ nodes is at most $2\log(n+1)$.
\end{lemma}
\begin{proof}
    We introduce $bh(x)$ the number of black nodes in a path from $x$ to a leaf.\\
    By induction on $bh(x)$, the subtree rooted at $x$ contains at least $2^{bh(x)} - 1$ :
    \begin{itemize}
        \item Initialization : $bh(x) = 0$ implies the tree has more than $0$ nodes which is true.
        \item Heredity :
              \begin{itemize}
                  \item If $x$ is black, let $y_{1}, y_{2}$ denote its children. We have $bh(y_{1}) = bh(x) - 1 = bh(y_{2}) - 1$. By summing, we get the result.
                  \item If $x$ is red, both of its children are black by rule 4. Let us denote the grand children of $x$ by $z_{1}, \ldots, z_{4}$. We have : $bh(z_{i}) = bh(x) - 1$. By sum, we get the result.
              \end{itemize}
    \end{itemize}
    Then if $h$ the black height of the tree, $n\geq 2^{h}-1$ and $h \leq \log{n+1}$. Since in any root-to-leaf path, the number of nodes is at most twice the number of black nodes per rule 4, the lemma follows.
\end{proof}

\begin{proposition}[Insertion]
    We insert a node into a n-node red-black tree by :
    \begin{itemize}
        \item Inserting the node into $T$ as if it were an ordinary binary search tree
        \item We color it red
        \item We perform a number of rotations and node recolouring.
    \end{itemize}
    This is done in $\O(\log n)$.
\end{proposition}

We define the left rotation on a node $x$ whose right child is not null like this :
\[Tree(x , \alpha, Tree(y, \beta, \gamma)) \mapsto Tree(y, Tree(x, \alpha, \beta), \gamma)\]
We define right rotations to be the reverse operation.
\begin{proposition}
    Rotations do not break the BST properties.
\end{proposition}
\begin{proof}[Only on right rotations]
    Elements in $\gamma$ are greater than $y$, in $\beta$ are greater than $x$ and lower or equal than $y$, and elements in $\alpha$ are lower or equal than $x$. After the rotation, we still have those properties.
\end{proof}

\begin{proposition}[Restoring red-black properties]
    Rule $1$ is never violated. Rules $2$ and $3$ can be fixed in $\O(1)$ time. Rule $5$ is never violated as we colour the new node red.\\
    We will try and maintain the invariant that the properties are broken in at most one node. There are three cases in which Rule $4$ can be violated : Let $z$ be the node that violates the red-black properties.
    \begin{enumerate}
        \item $z$'s uncle $y$ is red $\Rightarrow$ We colour $z.p$ black, colour $y$ black and $z.p.p$ red. We moved the problematic node up by two layers.
        \item $z$'s uncle $y$ is black and $z$ is a right child $\Rightarrow$ We left-rotate around the edge $z.p$-$z$, we then have $z$'s uncle $y$ to be black and $z$ to be a left child, i.e. case 3.
        \item $z$'s uncle $y$ is black and $z$ is a left child $\Rightarrow$ We color $z.p$ black, colour $z.p.p$ red and right-rotate.
    \end{enumerate}
\end{proposition}

\begin{proposition}[Complexity]
    Red black tree use $\O(n)$ space, $\O(\log n)$ in time for insertion and deletion in the worst case. % TODO : add deletion.
\end{proposition}

\subsection{Treaps}
\begin{definition}
    Treaps are binary trees with every node having a key and a priority verifying :
    \begin{itemize}
        \item a binary search tree for the keys
        \item a min-heap for the priorities
    \end{itemize}
\end{definition}

\begin{proposition}[Search]
    The usual algorithm yields time in $\O(depth of the node)$ for a successful search and $\O(\max \left(\texttt{depth}(v^{-}), \texttt{depth}(v^{+})\right))$ where $\texttt{key}(v^{-})$ is the predecessor of the searched key.
\end{proposition}

\begin{proposition}[Insertion]
    The standard BST algorithm with added rotations to fix the heap properties (if $p(z) \leq p(z.p)$, rotate auround $(z, z.p)$) yields time complexity in $\O(\max\left(\texttt{depth}(v^{-}), \texttt{depth}(v^{+})\right))$.
\end{proposition}

\begin{proposition}[Deletion]
    We run the insertion algorithm backwards : let $\texttt{light}(z)$ be the child of $z$ with smaller priority. The choice of this edge preserves the heap property everywhere except at $z$. As long as $z$ is not a leaft, rotate around $(z, \texttt{ligth}(z))$, then chop it off.
\end{proposition}

\begin{proposition}[Split]
    If we want to split the tree along a pivot $\pi$, we can insert a new node $z$ with key $\pi$ and priority $-\infty$. After the insertion, this new node is the root since it has the smallest priority, and all keys in the left subtree are smaller than the pivot and this subtree is a treap.
\end{proposition}

\begin{proposition}[Merge]
    Suppose that we want to merge two treaps $T_{-}$ and $T_{+}$, where keys in the first tree are smaller than the keys in the second tree. We create a dummy root with priority $-\infty$ and then delete it.
\end{proposition}

\begin{proposition}[Complexity]
    \begin{itemize}
        \item Time for a successful search : $\O(\texttt{depth(v)})$ where $\texttt{key}(v) = k$.
        \item Time for an unsuccessful search : $\O(\max \left(\texttt{depth}(v^{-}), \texttt{depth}(v^{+})\right))$ where $\texttt{key}(v^{-})$ is the predecessor of the searched key.
        \item Insertion Time : $\O(\max \left(\texttt{depth}(v^{-}), \texttt{depth}(v^{+})\right))$ where $\texttt{key}(v^{-})$ is the predecessor of the searched key.
        \item Deletion Time : $\O(\text{tree depth})$
        \item Split/Merge Time : same as insertion / deletion
    \end{itemize}
\end{proposition}

We will choose the priorities to be independently and uniformly distributed continuous random variables. We denote by $x_{1}, \ldots, x_{n}$ the nodes of the tree in the increasing order of the keys.
We denote  by $i \uparrow k$, the event $x_{i}$ if a proper ancestor of $x_{k}$ : $\texttt{depth}(x_{k}) = \sum_{i = 1}^{i = n} \left[i \uparrow k\right]$

\begin{lemma}
    Define \[
        X(i, k) = \left\{\begin{aligned}  & x_{i},\ldots, x_{k} & \text{ if } i < k  \\
                 & x_{k},\ldots, x_{i} & \text{ otherwise}.
        \end{aligned}\right.\]

    For all $i \neq k$ we have $i \uparrow k$ iff $x_{i}$ has the smallest priority in $X(i, k)$.
\end{lemma}
\begin{proof}
    \begin{enumerate}
        \item If $x_{i}$ is the root, then $x_{i}$ has the smallest priority
        \item If $x_{k}$ is the root, then $x_{i}$ is not an ancestor of $x_{k}$ and does not have the smallest priority
        \item If $x_{j}$ is the root with $i < j < k$, then $x_{i}$ is not an ancestor of $x_{k}$ and does not have the smallest priority.
        \item If $x_{j}$ is the root and either $i < k < j$ or $i < j < k$ then $x_{i}$ and $x_{j}$ are in the same subtree and the claim follow by induction.
    \end{enumerate}
\end{proof}

\begin{corollary}
    \[\mathbb{P}\left[i \uparrow k\right] =
        \left\{
        \begin{aligned}
            \frac{1}{k - i + 1} & \text{ if } i < k \\
            0                   & \text{ if } i = k \\
            \frac{1}{i - k + 1} & \text{ if } i > k
        \end{aligned}
        \right.
    \]
    Then :
    \[
        \mathbb{E}\left[\texttt{depth}(x_{k})\right] = \sum_{i = 1}^{i = n}\mathbb{P}\left[i \uparrow k\right] = \sum_{i = 1}^{k - 1}\frac{1}{k - i + 1} + \sum_{k+1}^{n}\frac{1}{i - k + 1} = H_{k} - 1 + H_{n - k + 1} - 1 < 2 \ln n - 2
    \]
    Thus, all treap operations take $\O(\log n)$ time in expectation
\end{corollary}

\begin{remark}
    The fastest sorting algorithm is Quicksort, and it can be seen as inputing all the values in the list in a treap then taking the infix order of the treap.
\end{remark}

\begin{remark}
    Tango trees are $\O(\log \log n)$-competitive and Splay trees are conjecture to be fastest in $\O(1)$-competitive, but we do not know if this is true.
\end{remark}

\section{Lower Bound for Sorting}
The question here is : can we sort in $\O(n\log n)$ or better ?
If comparisons only are allowed then no : indeed, any comparison-based sorting program can be thought of as defining a decision tree of possible executions. Running the same program twice on the same permutation causes it to do the exact same thing bu running it on different permutations of the same data causes a different sequence of comparisons to be made on each.

\subsection{Decision Tree}
\begin{definition}
    A decision tree for a sorting algorithm is a binary tree that shows the possible executions of an algorithm on a set.
\end{definition}

\begin{proposition}
    The minimum height of a decision tree is the worst-case complexity of comparison-based sorting.
\end{proposition}
\begin{lemma}
    The height of any decision tree is $\Omega(n \log n)$
\end{lemma}
\begin{proof}
    Since any two different permutations of $n$ elements require a different sequence of steps to sort, there must be at least $n!$ different paths from the root to leaves in the decision tree.
    Thus there must be at least $n!$ different leaves in this binary tree. Since a binary tree of height $h$ has at most $2^{h}$ leaves, we know that : $n! \leq 2^{h}$ i.e. $h\geq \log(n!)$. Finally, we have $\log n! = \Omega (n \log n)$.
\end{proof}
\begin{theorem}
    Any comparison-based sorting algorithm must use $\Omega(n\log n)$ time.
\end{theorem}

\section{Predecessor Problem}
We want to maintain a set $S$ of integers from a universe $U = \left[1, u\right]$ subject to insertions, deletions, predecessor and succesor queries. This problem is harder than dictionaries and hashing. BTSs give a solution in $\O(n)$ space and $\Theta(\log n)$ time.

\subsection{van Emde Boas Trees}
If the time per operation satisfies : $T(u) = T(\sqrt{u}) + \O(1)$, by Substitution, $T(u) = \O(\log \log u)$. We will split $U$ into $\sqrt{u}$ chunks of size $\sqrt{u}$ size, and the recurse.
\begin{definition}[Recursive van Emde Boas Trees]
    \begin{itemize}
        \item \texttt{T.summary} is a vEB-tree of size $\sqrt{u}$ containing all $i$ such that the $i$-th chunk is not empty.
        \item For each $1 \leq i \leq \sqrt{u}$, \texttt{T.chunk[$i$]} is a vEB-tree of size $\sqrt{u}$ containing $x \texttt{mod} \sqrt{u}$ for each $x$ in the $i$-th tree.
        \item \texttt{T.min} is the smallest element in $T$, not stored recursively.
    \end{itemize}
    We represent each integer $x = \langle c, i \rangle$ where
    \begin{itemize}
        \item $c$ is the chunk coordinates : $c = x // \sqrt{u}$
        \item $i$ is the position of $x$ in the chunk : $i = x \texttt{mod} \sqrt{u}$
    \end{itemize}
\end{definition}

\begin{proposition}[Successor Operation]
    With the following algorithm we get the wanted complexity :
    \begin{algorithm}
        \caption{Successor Complexity Verifies : $T(u) = T(\sqrt{u}) + \O(1)$}
        \begin{algorithmic}
            \Function \texttt{Successor}
            \Input $\ \ (T, x = \langle c, i \rangle)$
            \EndInput
            \If {$x < T.min$}
            \State \Return $T.min$
            \EndIf
            \If {$i < T.chunk[c].max$}
            \State \Return $\langle c, \texttt{Successor}\left(T.chunk[c], i\right) \rangle$
            \Else
            \State $c^{'} = \texttt{Successor}(T.summary, c)$
            \State \Return $\langle c^{'}, T.chunk[c^{'}].min\rangle$
            \EndIf
            \EndFunction
        \end{algorithmic}
    \end{algorithm}
\end{proposition}

\begin{proposition}[Insertion]
    The algorithm below gets the correct complexity.
    \begin{algorithm}
        \caption{Insertion Complexity Verifies : $T(u) = T(\sqrt{u}) + \O(1)$}
        \begin{algorithmic}
            \Function \texttt{Insert}
            \Input $\ \ (T, x = \langle c, i \rangle)$
            \EndInput
            \If {$T.min = None$}
            \State $T.min = T.max = x$
            \State \Return
            \EndIf
            \If {$x < T.min$}
            \State $\texttt{swap}(x, T.min)$
            \EndIf
            \If {$x > T.max$}
            \State $T.max = x$
            \EndIf
            \If {$T.chunk[c].min = None$}
            \State $\texttt{Insert}(T.summary, c)$
            \EndIf
            \State $\texttt{Insert}(T.chunk[c], i)$
            \EndFunction
        \end{algorithmic}
    \end{algorithm}
\end{proposition}

\begin{proposition}[Deletion]
    The algorithm below gets the correct complextiy.
    \begin{algorithm}
        \caption{Deletion Complexity Verifies : $T(u) = T(\sqrt{u}) + \O(1)$}
        \begin{algorithmic}
            \Function \texttt{Delete}
            \Input $\ \ (T, x = \langle c, i \rangle)$
            \EndInput
            \If {$x = T.min$}
            \State $c = T.summary.min$
            \If {c = None}
            \State $T.min = None$. \Return
            \EndIf
            \State $x = T.min = \langle c, i = T.chunk[c].min \rangle$
            \EndIf
            \State $\texttt{Delete}(T.chunk[c], i)$
            \If {$T.chunk[c].min = None$}
            \State $\texttt{Delete}(T.summary, c)$
            \EndIf
            \If {$T.summary.min = None$}
            \State $T.max = T.min$
            \Else
            \State {$c^{'} = T.summary.max$}
            \State {$T.max = \langle c^{'}, T.chunk[c^{'}].max \rangle$}
            \EndIf
            \EndFunction
        \end{algorithmic}
    \end{algorithm}
\end{proposition}

\begin{proposition}[Complexity]
    \begin{itemize}
        \item All of these operations' complexities verify $T(u) = T(\sqrt{u}) + \O(1)$ and thus have time complexity $\O(\log \log u)$
        \item Space complexity satifies : $S(u) = \left(\sqrt{u} + 1\right)\dot S(\sqrt{u}) + \O(1)$ and therefore $S(u) = \O(u)$.
    \end{itemize}
\end{proposition}

\begin{proposition}[Original van Emde Boad trees]
    For $\texttt{Successor}(x)$, as the path from the root to $x$ is monotone, binary searching the path to find the lowest 1 gives us either the predecessor or the successor of $x$. By storing all nodes in an array of size $\O(u)$ to allow efficient binary search; a pointer from each node to the maximum and minimum of their subtree; all the elements as a doubly-linked list, we find the successor and the predecessor of $x$ in $\O(\log \log u)$ time.\\
    Update is done in $\O(\log u)$ time, since we only need to update the element-to-root path, in $\Theta(u)$ space
\end{proposition}

\subsection{Improvements}
\subsubsection{x-fast trees}
\begin{definition}
    In an x-fast tree, we store every root-to-green node (nodes representing an element from the set) path, viewed in binary (left = 0, right = 1), via Cuckoo Hashing.
\end{definition}
\begin{proposition}
    Predecessor queries are done in $\O(\log \log u)$ time, updates in $\O(\log u)$ expected amortised time, but this tree only takes $\O(n \log u)$ space.
\end{proposition}
\begin{proof}
    To maintain successor and predecessor, we use the vEB-tree algorithm, giving us the same complexity, and same for updates. Yet, since we only store the root-to-element paths which are of $\log u$ length, we need $\O(n\log u)$ space.
\end{proof}

\subsubsection{y-fast trees}
\begin{definition}
    We maintain elements in groups of size in $\left[\frac{\log u}{4}, 2\log u\right]$. For each group, we build a BST, and we store representatives of the group using an x-fast tree :
    \begin{itemize}
        \item If there are fewer than $\frac{\log u}{2}$ elements, we store them in a single BST.
        \item otherwise, suppose we add/delete an element. If a group becomes too large, we split it in two. If a group becomes too small, we merge it with its neighbour, then split if needed.
    \end{itemize}
\end{definition}

\begin{proposition}
    Predecessor queries are in $\O(\log \log u)$ time, updates in $\O(\log \log u)$ expected amortised time, since insertion into the x-fast trie happens only once per $\Theta(\log u)$ new elements.
\end{proposition}
\begin{proof}
    This comes directly from the definition and the definition of the x-fast trees.
\end{proof}

\part[Algorithms on strings]{Lecture 5 : 9/11}
\localtableofcontents
\section{General Definitions}
\begin{definition}
    An alphabet $\Sigma$ is a finite set. Elements of $\Sigma$ are letters, characters or symbols, denoted by small letters. A string (or word), is a finite sequence of letters, denoted by a capital letter. \\
    The length of a string $S = s_{1}s_{2}\ldots s_{n}$ where $s_{i} \in \Sigma$ is denoted by $\abs{S}$ and defined to be equal to $n$.\\
    For $1 \leq i \leq j \leq S$, we say that $S[i, j] = s_{i}\ldots s_{j}$ is a substring of $S$. If $i = 1$, $S[i, j]$ is a prefix of $s$. If $j = \abs{S}$, $S[i, j]$ is a suffix of $S$.
    $S[i,j]$ is an occurence of a string $X$ in $S$ if $S[i, j] = X$
\end{definition}

\section{Pattern Matching}
\begin{definition}
    We define the pattern matching problem as such : Given a string $P$ of length $m$ (pattern), a string $T$ of lenth $n \geq m$ (text), return all occurences of $P$ in $T$ specified by their starting positions.
\end{definition}

\subsection{Naïve Algorithm}
The naïve algorithm consists of looking, for each possible starting position in the text, check if the $m$ following characters form a string equal to the pattern. This algorithm has time complexity $\O(nm)$ and space complexity $\O(n + m)$.
\begin{remark}
    There are more than 80 algorithms for this problem. See S.Faro, T.Lecroq \href{The Exact String Matching Problem}{https://arxiv.org/pdf/1012.2547.pdf}
\end{remark}

\subsection{Knuth-Morris-Pratt Algorithms}
\begin{definition}
    The border of a string $P$ is a proper prefix of $P$ that is equal to a suffix of $P$.\\
    The border array of length $\abs{P}$ is such that $B[i]$ is the maximal length of a border of $P[1, i]$, $0$ if undefined.
\end{definition}

\begin{proposition}\label{lemma1KMP}
    $B$ can be computed in $\O(\abs{P})$ time.
\end{proposition}
\begin{lemma}
    First, note that :
    \begin{enumerate}
        \item $B[1] = 0$
        \item If $X$ is a border of $P[1, k]$ of length $x$, then $X^{'} = X[1, x - 1]$ is a border of $P[1, k - 1]$.
        \item $B[k] - 1 = B[k - 1]$ if and only if $P[k] = P[B[k - 1] + 1]$.
        \item If $X$ is a border of $X$ and $Y$ is a border of $X$, $Y$ is a border of $S$.
    \end{enumerate}
    Suppose that we have computed $B[1], \ldots, B[k - 1]$. We will now compute $B[k]$.\\
    By property 3, if $P[k] = P[B[k - 1] + 1]$, then $B[k] = B[k - 1] + 1$.\\
    Else, if $P[k] \neq P[B[k - 1] + 1]$, consider $B^{2}[k - 1] = B[B[k - 1]]$. If $P[k] = P[B^{2}[k - 1] + 1]$, set $B[k] = B^{2}[k - 1] + 1$, else consider $B^{3}[k - 1]$, and so on.\\

    This algorithm is correct and in $\O(m)$.
\end{lemma}
\begin{proof}
    \begin{itemize}
        \item Correctness : In fact, $B[1, k] \in \left\{B[k - 1] + 1, B^{2}[k - 1] + 1, \ldots, 0\right\}$.\\
              From property 2, any border of $P[1, k]$ can be obtained by appending $P[k]$ to some border of $P[1, k - 1]$.\\
              From property 4, $B^{j}[k - 1]$ is the $j$-th longest border of $P[1, k - 1]$.
        \item Time : $\texttt{Time}(B[k]) \leq \abs{B[k] - B[k - 1]}$, indeed, if $B[k] = B^{j}[k - 1] + 1$, we use $j$ steps.
              Observe that :
              \[
                  -1 \leq \sum_{k}\left(B[k] - B[k - 1]\right) = \sum^{-}\left(B[k] - B[k - 1]\right) + \sum^{+}\left(B[k] - B[k - 1]\right) \leq m
              \]
              However, by property $2$, the sum on the positive differences if $\O(m)$. Hence, the total time is in $\O\left(\sum \abs{B[k] - B[k - 1]}\right) = \O(m)$
    \end{itemize}
\end{proof}

% TODO : add match function 
\begin{algorithm}
    \caption{Knuth-Morris-Pratt Algorithm : KMP}
    \begin{algorithmic}
        \Input $\ \ T, P, n, m$
        \EndInput
        \State $B \gets \texttt{BorderArray}(P)$
        \State $L \gets []$
        \State $i \gets 1, \ j\gets 1$
        \While {$i \leq n - m + j$}
        \State $(i, j) \gets \texttt{match}(i, j, m)$
        \If {$j = m + 1$}
        \State $L \gets (i - m)::L$
        \Else
        \If {$j = 1$}
        \State $i \gets i + 1$
        \Else
        \State $j \gets B[j - 1] + 1$
        \EndIf
        \EndIf
        \EndWhile
        \Return \ $L$
    \end{algorithmic}
\end{algorithm}

\begin{proposition}
    This algorithm is Correct
\end{proposition}
\begin{lemma}\label{lemma2KMP}
    No occurence of $P$ can start in the shift interval.
\end{lemma}
\begin{proof}
    If $T[x, x + m - 1] = P$, $P[1, j - 1]$ has a border of length $> B[j - 1]$, contradiction.
\end{proof}

\begin{lemma}\label{obs1KMP}
    We have $ P[1, B[j-1]] = T[i + (j - 1 - B[j-1]), i + j - 2]$. \\
    Hence, we do not need to compare the first $B[j - 1]$ letters of the patter with the letters of the text after the shift.
\end{lemma}
\begin{proof}[Of the Correctness]
    This stands from \ref{lemma2KMP} and \ref{obs1KMP}.
\end{proof}

\begin{proposition}
    This is done using $\O(m)$ extra space, in $\O(n + m)$ time.
\end{proposition}
\begin{proof}
    \begin{itemize}
        \item Space : $\O(m)$ to store $B$ + constant
        \item Time : From \ref{lemma1KSP}, $\O(m)$ to build $B$. To process $T$ we need $\O(n + m)$ time, since $1 \leq i \leq n$ and $1 \leq j \leq n$ and their difference increases after any step.
    \end{itemize}
\end{proof}

\begin{definition}
    An algorithm is called online if can process its input item-by-item in a serial fashion, without having the entire input available from the start.\\
    An online algorithm is called real-time if it processes each data item in constant time.\\
    KMP is an online algorithm, but it is not real-time.
\end{definition}

\begin{proposition}
    We can construct a real-time version of KMP, by using $B^{'}[j, T[i]]$ to decide how to shift the pattern, where :
    \[
        B^{'}[j, a] = \max{\left\{k \mid k < j,\ P[1, k] = P[j - k + 1, j], \ P[k + 1] = a\right\}}
    \]
    $B^{'}$ occupies $\O(m)$ space and is computed in $\O\left(m \cdot \abs{\Sigma}\right)$ time.
    The modified KMP is real-time and correct.
\end{proposition}

\section{Dictionaries, Multiple Pattern Matching}
\subsection{Tries}
\begin{definition}
    A dictionary $D$ is a set of string. A trie of $D$ is a tree with every node being labeled by a letter such that :
    \begin{itemize}
        \item For every node, outgoing edges are labeled with different letters
        \item For every string $S \in D$, there is a root-to-node path that spells out $S$. The end of the path is labeled with the id of $S$.
        \item Every root-to-leaf path spells out a string in $D$.
    \end{itemize}
\end{definition}
\begin{proposition}
    A trie for $D$ uses $\O(\sum_{S \in D}\abs{S})$ space.
\end{proposition}
\begin{proof}
    By laying out each root-to-leaf path.
\end{proof}

We will consider two applications :
\begin{enumerate}
    \item Dictionary Look-Up : Given a string $P$, decide if it belongs to the dictionary.
    \item Multiple Pattern Matching : Given a dictionary of patterns $D$, find all their occurences in a text $T$.
\end{enumerate}

\subsection{Dictionary Look-Up}
\begin{proposition}
    For $P$ a pattern, $P \in D \Leftrightarrow $  there is a root-to-node path that spells out $P$ labeled by an id of a string from $D$. By starting at the root, if there is an from the root to its child $u$ labeled by $P[1]$, the algorithm moves to $u$ and continues recursively. This takes time in $\O(\abs{P})$.
\end{proposition}

\subsection{Aho-Corasick Algorithm}
We will first assume that no pattern in $D$ is a substring of another.
\begin{definition}
    The failure link from a node $u$ labeled by a string $S$ points to the deepest node $v$ labeled by a proper suffix of $S$.
\end{definition}

\begin{lemma}
    The trie occupies $\O(m)$ space and can be constructed in $\O(m)$ time.
\end{lemma}
\begin{proof}
    \begin{itemize}
        \item Space : There are $\O(m)$ nodes, each node has exactly one failure link.
        \item Time : Failure links are built top-to-down. Links from nodes of depth $1$ point to the root.
              Suppose that we have built links for all nodes to depth $\leq d - 1$. The failure link from a node labeled with $Sa$ must point to a node labeled with $S^{'}a$, where $S^{'}$ is a suffix of $S$. In other words, the parent of the end of the failure link from the node must belong to the failure link path from the node $p(u)$ , and it must be the deepest node that has an outgoing edge labeled with $a$.
              The algorithm works as follows: we follow the failure link path from $p(u)$. The first node with an outgoing edge labeled with $a$ is the end of the failure link for $u$.
              Consider the time needed to build the failure links for one root-to-leaf path. Let $root, u_{1}, \ldots, u_{k}$ be the nodes in this path, and denote by $f(u_{i})$ the depth of the end of the failure link for $u_{i}$. The time to find the link for $u_{i}$ is $\leq c \times \left(2 + f(u_{i-1} - f(u_{i}))\right)$. Thus the total time for the trie is $\O(m)$.
    \end{itemize}
\end{proof}

\begin{algorithm}
    \caption{Aho-Corasick Algorithm}
    \begin{algorithmic}
        \Input \ $T, D, n$
        \EndInput
        \State $curr\_node \gets root, i \gets 1, L \gets []$
        \While {$i \leq n$}
        \If {$\exists e = (curr \_node, u)$ labeled with $T[i]$}
        \State $curr\_node \gets u$
        \If {$curr\_node$ corresponds to a pattern}
        \State $L \gets i :: L$
        \EndIf
        \State $i \gets i + 1$
        \Else
        \If {$curr \_node = root$}
        \State $i \gets i + 1$
        \Else
        \State $curr\_node = \texttt{failure\_link}(curr\_node)$
        \EndIf
        \EndIf
        \EndWhile
    \end{algorithmic}
\end{algorithm}

\begin{proposition}
    This algorithm has space complexity $\O(m)$ and time complexity $\O(m + n)$.
\end{proposition}
\begin{proof}
    \begin{itemize}
        \item Space : $m$ is the total length of the patterns
        \item Time : We need $\O(m)$ time to build the trie. Consider how the depth of $curr\_node$ changes during the algorithm.
              \begin{itemize}
                  \item Every time we do down an edge (happens $\leq n$ times), it increases by $1$. Every time we follow a failure link, it decreases by $\geq 1$.
                  \item Therefore, as the depth is always positive, we follow a failure link at most $n$ times.
              \end{itemize}
    \end{itemize}
\end{proof}

\begin{theorem}
    If no pattern is a substring of another, the multiple pattern matching problem can be solved in $\O(m)$ space and $\O(n + m)$ time.
\end{theorem}

\begin{definition}
    When patterns are substrings of others, we add output links : an output link from a node $u$ goes to the nearest node on the failure link path outgoing from $u$ that corresponds to a pattern of the dictionary.
\end{definition}

\begin{proposition}
    Output links can be built in $\O(m)$ time by one top-down traversal of the failure link tree.
\end{proposition}

We modify the algorithm : at each step it follows the path from $curr \_node$ and outputs the patterns in the output link path.

\begin{theorem}
    The multiple pattern matching problem can ve solved in $\O(m)$ space and $\O(n + m + occ)$ where $occ$ is the total number of occurences of the patterns in the text.
\end{theorem}

\section{Suffix Trees and Applications}
\begin{definition}
    A text index is a data structure that represents a text and supports pattern matching queries. The suffix tree is one of them.

    We define the compact trie for $D$ as the trie for $D$ where nodes with only one child have been replaced by an edge. The suffix tree of a string is the compact trie for the set of the suffixes of that string with \$ appended (\$ must not be in $\Sigma$).
\end{definition}

\subsection{Suffix Trees}
\begin{proposition}
    A suffix tree for a string of length $n$ has $n$ leaves, and at most $2n - 1$ nodes and $2n - 2$ edges. Storing the labels on the edges can take $\Theta\left(\abs{T}^{2}\right)$. To save space we represent each label as two numbers, the left and the right endpoints in $T$.
\end{proposition}
\begin{proof}
    By induction.
\end{proof}

To implement algorithms on the suffix tree efficiently, we need to be able to identify, given a node $u$ and a letter $a$, an edge $u, v$ such that its label starts with $a$. Then, for each node $u$, we store an array $A_{u}$ of size $\abs{\Sigma}$ , with $A_{u}[a]$ being a pointer to the child $v$ of $u$ such that the label of $(u, v)$ starts with $a$.

\begin{proposition}
    This structure takes $\O(\abs{Sigma} \cdot \abs{T})$ space.
\end{proposition}

\subsection{Pattern Matching}
\begin{remark}
    A suffix of $T$ starts with $P$ if and only if it corresponds to a leaf of the suffix tree that belongs to a subtree rooted at the end of the path labeled with $P$.
\end{remark}

To answer pattern matching queries, we start at the root. Then, we follow the path labeled by the letters of $P$, and use arrays $A_{u}$ to find the next edge to follow. If there is no path labeled by $P$, there are no occurences of $P$ in $T$. Otherwise, the starting positions of the occurences of $P$ in $T$ are the starting positions of the suffixes that are in the subtree rooted at the end of the path labeled by $P$. We retrieve the occurences by traversing the subtree depth-first.
\begin{proposition}
    This algorithm has total time $\O(m + occ)$.
\end{proposition}
\begin{proof}
    The arrays allow to decide which edge to follow next in $\O(1)$ time. In total, finding the path requires $\O(m)$ time. If there is a path labeled by $P$, the subtree rooted at its end has $\O(occ)$ leaves, where $occ$ is the number of occurences of $P$ in $T$. As the subtree does not have nodes of degree one (except possibly for the root), its size is $\O(occ)$.
\end{proof}

\subsection{Longest Common Substring}
Given two strings, find the longest substring that occurs both in $T_{1}$ and $T_{2}$. We will denote $n = \abs{T_{1}} + \abs{T_{2}}$.

To find the longest common substring of $T_{1}$ and $T_{2}$, we need the suffix tree containing the suffixes of both string (sometimes called generalised suffix tree).

\begin{proposition}
    This suffix tree can be built in $\O\left(\abs{T_{1}} + \abs{T_{2}}\right)$.
\end{proposition}

\begin{remark}
    If a string $X$ occurs at position $i$ in $T_{1}$, and at position $j$ in $T_{2}$, then the subtree rooted at the end of the path labeled by $X$ contains the leaf corresponding to $T_{1}[i\ldots]$ and $T_{2}[j \ldots]$.
\end{remark}

By bottom-up traversal, we find the deepest node in the tree such that its subtree contains both leaves corresponding to suffixes of $T_{1}$ and suffixes of $T_{2}$. It is labeled by the longest common substring of $T_{1}$ and $T_{2}$.
\begin{proposition}
    This algorithm has time complexity $\O(\abs{T_{1}} + \abs{T_{2}})$.
\end{proposition}

We can generalise this to any $m$ strings in time $\O(n)$ where $n = \sum_{k = 1}^{m} \abs{T_{k}}$. We will use the following fact : we can preprocess a tree of size $\O(n)$ in time $\O(n)$ to support lowest common ancestor ($LCA$) queries in constant time. (see ) $LCA(u, v)$ must return the lowest node that is an ancestor of both $u$ and $v$.

\begin{openpb}
    We showed that the longest common substring problem for two string can be solved in $\O(n)$ space and $\O(n)$ time. It is also known that the problem can be solved in $\O(s)$ extra space and time
    \[
        \O\left(\frac{n^{2}\log(n)\log^{*}(n)}{L\cdot s} + n \log n\right)
    \]
    where $L$ is answer.
    What is the optimal trade-off.
\end{openpb}

\part[Disjoint-Set Structure]{Lecture 6 : 23/11}
\localtableofcontents
\section{Disjoint-set Data Structure}
\begin{definition}
    A disjoint-set data structure is a way to maintain a collection of disjoint sets with 3 operations :
    \begin{itemize}
        \item \texttt{make\_set($x$)} : create a new set containing one element, $x$
        \item \texttt{union\_set($x, y$)} : union the sets containing $x$ and $y$
        \item \texttt{find\_set($x$)} : return a pointer to the representative of the set containing $x$
    \end{itemize}
    We will parametrize by $n$ the number of elements and $m$ the number of operations.
\end{definition}

It can be used to preprocess the connected components of a graph, so as to maintain the queries of the paths between two vertices.\\

It can also be used to represent the ways player play in the HEX game.


\section{Linked-list Representation}
Here, each set is a linked list, and the representative is the head of the list. For every element, we store a pointer to the node containing it. In that node, we store a pointer to the representative.
\begin{proposition}
    In this situation : \texttt{make\_set} and \texttt{find\_set} are done in $\O(1)$.
\end{proposition}

For \texttt{union\_set}, it requires more work.\\
Indeed, simply appending the set of $y$ to the set of $x$ takes constant time, but updating the pointers to the representatives is done in $\Theta\left(\abs{\text{set of }y}\right)$. By doing the following sequence of operation :
\[
    \texttt{make\_set}(x_{1}), \texttt{make\_set}(x_{2}), \ldots, \texttt{make\_set}(x_{n});\\
    \texttt{union\_set}(x_{2}, x_{1}), \texttt{union\_set}(x_{3}, x_{2}),\ldots, \texttt{union\_set}(x_{n}, x_{n-1})
\]
we see that even the amortised time per operation can be large, as it takes a total time of $\Theta\left(n^{2}\right)$ and the amortised being $\Theta(n)$.

Then, we propose a weighted-union strategy : if the set of $x$ is larger than the set of $y$, we append the latter to the former, else we append the set of $x$ to the set of $y$.

\begin{theorem}
    Using linked-lists and the weighted-union strategy, doing $m$ \texttt{make\_set}, \texttt{union\_set} and \texttt{find\_set}, $n$ of which are \texttt{union\_set} takes $\O\left(m + n\log n\right)$
\end{theorem}

\begin{proof}
    At any moment, the largest size of a set is bounded by $n$. Therefore, an element can change its set representative at most $\log_{2} n$, as every time it happends the size of the set at least doubles in size. Therefore, the total time for updating the representatives is $\O(n \log n)$. All other operations take $\O(m)$ time.
\end{proof}

\section{Disjoint-set forests}
Here, each set is a treee whose root is the representative. Each node stores a pointer to its parent. \texttt{make\_set}($x$) creates a tree containing only one node, \texttt{find\_set}($x$) follows the path from $x$ to the root and \texttt{union\_set}($x, y$) makes the root of the tree of $y$ become a child of the root of the tree of $x$.\\
We will use two strategies :
\begin{itemize}
    \item Path Compression : During \texttt{find\_set}($x$), we make each node on the path to the root be a child of the root.
    \item Union by Rank : During \texttt{union\_set}, the root of the tree with smaller rank becomes a child of the root of the tree with larger rank.
\end{itemize}

\begin{definition}
    We define ranks as upper bounds on the height of the node :
    \begin{itemize}
        \item When we add a new node $x$, we initialise the rank of $x$ with $0$.
        \item If we union two trees (i.e.,two sets) of different ranks, ranks do not change. If we union two trees of equal ranks, the rank of the root of the resulting tree increases by one.
    \end{itemize}
\end{definition}

\begin{algorithm}
    \caption{Union-Find}
    \begin{algorithmic}
        \Function \texttt{make\_set}
        \Input \ \ $x$
        \EndInput
        \State $x.p \gets x$
        \State $x.rank \gets 0$
        \EndFunction

        \Function \texttt{union\_set}
        \Input \ \ $x, y$
        \EndInput
        \State \texttt{link(find\_set($x$), find\_set($y$))}
        \EndFunction

        \Function \texttt{find\_set}
        \Input \ \ $x$
        \EndInput
        \If {$x \neq x.p$}
        \State $x.p \gets find\_set(x.p)$
        \EndIf
        \Return $x.p$
        \EndFunction

        \Function \texttt{link}
        \Input \ \ $x, y$
        \EndInput
        \If {$x \neq y$}
        \If {$x.rank > y.rank$}
        \State $y.p \gets x$
        \Else
        \State $x.p \gets y$
        \If {$x.rank = y.rank$}
        \State $y.rank \gets y.rank + 1$
        \EndIf
        \EndIf
        \EndIf
        \EndFunction
    \end{algorithmic}
\end{algorithm}


\section{Time Analysis}
We defined $n$ as the number of elements, i.e. of \texttt{make\_set} and $m$ the total number of operations.

\begin{theorem}[Tarjan] % FUCK TARJAN
    Using both Union by Rank and Path Compression, the time complexity for any sequence of Find and Union is $\O(m + n\alpha(n))$ where $\alpha(4)$ is the inverse of the Ackermann function. In any conceivable application $\alpha(n) \leq 4$. This time complexity is tight.
\end{theorem}

\subsection{A simpler Bound}
We will first start with a simpler analysis showing $\O(m \log^{\star} n)$.\\

\begin{proposition}
    Suppose we convert a sequence of $m^{'}$ operations \texttt{make\_set, union\_set, find\_set} into $m$ operations \texttt{make\_set, link, find\_set}. If we can execute the latter in $\O(m \cdot \alpha(n))$, we can execute the former in $\O(m^{'} \cdot \alpha(n))$.
\end{proposition}

\begin{proof}
    $m^{'} \leq m \leq 3m^{'}$
\end{proof}
We will thus assume we have only \texttt{make\_set, link, find\_set} operations.

\begin{proposition}
    \begin{itemize}
        \item The rank of a node can only increase with time. Rank of non-root nodes are frozen forever.
        \item For every node $x$, $x.rank \leq x.p.rank$ and the inequality is strict unless $x$ is a root node.
        \item The ranks of nodes in any path strictly increase
        \item The rank of any node is bounded by $n - 1$.
    \end{itemize}
\end{proposition}

\begin{proposition}[Rank Lemma]
    At any moment, for every $r \in \N$, there are at most $\frac{n}{2^{r}}$.
\end{proposition}

\begin{proof}
    By induction, the number of nodes in a tree with a root of rank $r$ is at least $2^{r}$ :
    \begin{itemize}
        \item Base Case : $r = 0$ is true
        \item Induction Step : After linking two trees with roots of rank $r_{1}, r_{2}$ :
              \begin{itemize}
                  \item If $r_{1} < r_{2}$ the claim is true by induction for $r_{2}$.
                  \item If $r = r_{1} = r_{2}$, then the root of the resulting tree has rank $r + 1$ and the tree has at least $2^{r} + 2^{r} = 2^{r + 1}$ nodes.
              \end{itemize}
    \end{itemize}
    Consider all nodes of rank $r$. None of them is not and was not an ancestor of another, so there are $2^{r}$ nodes corresponding to each of them in the forest, and these sets are independent. As we have $n$ nodes in total, the bound follows.
\end{proof}

\subsection{Buckets}
\begin{definition}
    We define buckets so they form a partition of $\N$ with : $B_{0} = \left[0, 0\right], B_{1} = \left[1, 1\right]$, $B_{2} = \left[2, 2^{2} - 1\right], B_{3} = \left[2^{2}, 2^{(2^{2})} - 1\right]$ and so on. They are not implemented, bu only used for the analysis. The total number of buckets is at most $\log^{\star} n$.
\end{definition}

\begin{proposition}\label{bucketlemma}
    Put a node with rank $r$ into the bucket $B$ containing $r$. The number of nodes in a bucket $B = \left[b, 2^{b} - 1\right]$ is at most :
    \[
        \sum_{r \in B} \frac{n}{2^{r}} \leq \frac{2n}{2^{b}}
    \]
\end{proposition}

\begin{definition}
    At a given time a node is good if :
    \begin{itemize}
        \item It is a root
        \item It is a child of a root
        \item It is in a smaller bucket than its parent
    \end{itemize}
    It is bad otherwise.
\end{definition}

A \texttt{find\_set} visits $\O(\log^{\star} n)$ good nodes. Therefore, the total time complexity is $\O(m\log^{\star} n + \abs{\left\{\text{bad nodes visited}\right\}})$.

\begin{proposition}
    The total number of visits of bad nodes is $\O(n\log^{\star} n)$
\end{proposition}
\begin{proof}
    Bad nodes have frozen rank. When a bad node $x$ is visited by a \texttt{find\_set}, the new $x.p$ has rank strictly larger than the previous one. Hence, for each bad node $x$ in a bucket $B = \left[b, 2^{b}-1\right]$ the number of times $x$ is visited while $x$ is bad is at most $2^{b} - 1$, so by \ref{bucketlemma}, it is at mode $\O\left(\frac{2n}{2^{b}}\right)$. So the number of times we visit a bad node in $B$ is $\O(n)$.
\end{proof}

Thus, the total complexity is $\O\left(m \log^{\star} n\right)$

In the following sections, we will try to improve this complexity.
\subsection{Ackermann's function}
\begin{definition}
    For $k \geq 0$, $j \geq 1$ two integers, the Ackermann function of $k$ and $j$ is defined as :
    \[
        A_{k}(j) = \begin{cases}
            j + 1                & \text{ if } k = 0    \\
            A_{k - 1}^{j + 1}(j) & \text{ if } k \geq 1
        \end{cases}
    \]
\end{definition}

\begin{proposition}[Properties]
    \begin{itemize}
        \item Lemma 1 : $\forall j \geq 1, A_{1}(j) = 2j + 1$
        \item Lemma 2 : $\forall j \geq 1, A_{2}(j) = 2^{j + 1}(j + 1) - 1$
    \end{itemize}
\end{proposition}
\begin{proof}
    Both are done by induction.
\end{proof}

\begin{definition}
    We define reverse Ackermann's function by :
    \[
        \alpha(n) = \min \left\{k \mid A_{k}(1) \geq n\right\}
    \]
    We have :
    \[
        \alpha(n) = \begin{cases}
            0 & \text{ if } 0 \leq n \leq 2 \\
            1 & n = 3                       \\
            2 & 4 \leq n \leq 7             \\
            3 & 8 \leq n \leq 2047          \\
            4 & 2048 \leq n \leq A_{4}(1)   \\
        \end{cases}
    \]
\end{definition}

\subsection{Potential Function}
We define a potential function :
\begin{definition}
    \begin{itemize}
        \item $\texttt{level}(x) = \max \left\{k \mid x.p.rank \geq A_{k}(x.rank)\right\}$
        \item $\texttt{iter}(x) = \max \left\{i \mid x.p.rank \geq A_{\texttt{level}(x)}^{(i)}(x.rank)\right\}$
    \end{itemize}
    For a node $x$, we define :
    \[
        \Phi_{q}(x) = \begin{cases}
            \alpha(n)\cdot x.rank                                                       & \text{ if $x$ if root or } x.rank = 0 \\
            \left(\alpha(n) - \texttt{level}(x)\right)  \cdot x.rank - \texttt{iter}(x) & \text{ otherwise}
        \end{cases}
    \]
    For a forest, define :
    \[
        \Phi_{0}(F) = 0 \text{ and } \Phi_{q}(F) = \sum_{x \in F} \Phi_{q}(x)
    \]
\end{definition}

We will show that $\Phi_{q}(F) \geq 0$ :
Then, if we denote by $t_{q}$ the time for the $q$-th operation :
\[
    \sum_{q} t_{q} \leq \sum_{q} t_{q} + \Phi_{q}(F) - \Phi_{0}(F) = \sum_{q}t_{q} + \sum_{q}\left(\Phi_{q + 1}(F) - \Phi_{q}(F)\right) = \sum_{q}\left(t_{q} + \Phi_{q + 1}(F) - \Phi_{q}(F)\right)
\]

\begin{lemma}
    If $x \neq x.p$ and $x.rank \geq 1$ : $0 \leq \texttt{level}(x) \leq \alpha(n) - 1$
\end{lemma}
\begin{proof}
    We have $x.p.rank \geq x.rank + 1 = A_{0}(x.rank)$, hence $\texttt{level}(x) \geq 0$.\\
    As $A_{\alpha(n)}(x.rank) \geq A_{\alpha(n)}(1) \geq n \geq x.p.rank$, we get the other inequality.
\end{proof}

\begin{lemma}
    If $x.rank \geq 1$, $1 \leq \texttt{iter}(x) \leq x.rank$.
\end{lemma}
\begin{proof}
    We have $x.p.rank \geq A_{\texttt{level}(x)}(x.rank) = A_{\texttt{level}(x)}^{1}$ thus $\texttt{iter}(x) \geq 1$.\\
    Since : $A_{\texttt{level}(x)}^{x.rank + 1}(x.rank) = A_{\texttt{level}(x) + 1}(x.rank) > x.p.rank$ and thus $\texttt{iter}(x) \leq x.rank$.
\end{proof}

\begin{lemma}
    $\forall x, q, 0 \leq \Phi_{q}(x) \leq \alpha(n)x.rank$
\end{lemma}
\begin{proof}
    If $x$ is the root or $x.rank = 0$, there is equality and we are done. Otherwise :
    \[
        \Phi_{q}(x) = \left(\alpha(n) - \texttt{level}(x)\right) \cdot x.rank - \texttt{iter}(x) \geq \left(\alpha(n) - (\alpha(n) - 1)\right) \cdot x.rank - \texttt{iter}(x) \geq x.rank - x.rank = 0
    \]
    Also :
    \[
        \Phi_{q}(x) \leq \alpha(n)\cdot x.rank - 1 < \alpha(n) x.rank
    \]
\end{proof}
\begin{corollary}\label{cor9}
    If $x$ is not a root and $x.rank > 0$, then $\Phi_{q}(x) < \alpha(n)x.rank$
\end{corollary}

\begin{lemma}\label{lemma10}
    If $x$ is not a root, the $q$-th operation is \texttt{link} or \texttt{find\_set} then $\Phi_{q}(x) \leq \Phi_{q - 1}(x)$. Moreover, if $x.rank \geq 1$ and either $\texttt{level}(x)$ or $\texttt{iter}(x)$ changes, then $\Phi_{q}(x) \leq \Phi_{q - 1}(x) - 1$.
\end{lemma}
\begin{proof}
    Since $x$ is not a root, the rank of $x$ is frozen.
    \begin{enumerate}
        \item If $x.rank = 0$ then $\Phi_{q}(x) = \Phi_{q - 1}(x) = 0$
        \item Assume $x.rank \geq 1$.
              \begin{itemize}
                  \item If \texttt{level}($x$) increases, $\left(\alpha(n) - \texttt{level}(x)\right)x.rank$ drops by at least $x.rank$ but $\texttt{iter}(x)$ can drop by at most $x.rank - 1$ therefore, $\Phi_{q}(x) \leq \Phi_{q - 1}(x) - 1$.
                  \item If \texttt{level}($x$) is unchanged, \texttt{iter}$(x)$ remains unchanged or increases. If it does not change $\Phi_{q}(x) = \Phi_{q - 1}(x)$, otherwise, $\Phi_{q}(x) \leq \Phi_{q - 1}(x) - 1$
              \end{itemize}
    \end{enumerate}
\end{proof}

\subsection{Time Bound}
More lemmas :
\begin{lemma}
    The amortised cost of each \texttt{make\_set} is $\O(1)$
\end{lemma}
\begin{proof}
    \texttt{make\_set} creates a node $x$ with $\Phi_{q}(x) = 0$. Hence, $\Phi_{q}(F) = \Phi_{q - 1}(F)$ and the amortised time is $1$.
\end{proof}

\begin{lemma}
    The amortised time of $\texttt{link}(x, y)$ is $\O(\alpha(n))$.
\end{lemma}
\begin{proof}
    WLOG, $y$ becomes the parent of $x$. By \ref{lemma10}, the potentials of any node but $x$ and $y$ cannot increase. Before \texttt{link}, the node $x$ is a root and $\Phi_{q - 1}(x) = \alpha(n)x.rank$. If $x.rank = 0$, we have $\Phi_{q - 1}(x) = \Phi_{q}(x) = 0$. Otherwise by \ref{cor9}, $\Phi_{q}(x) < \alpha(n)x.rank = \Phi_{q - 1}(x)$.\\
    For $y$, either $\Phi_{q - 1}(y) = \Phi_{q}(y)$ or $\Phi_{q}(y) = \Phi_{q - 1}(y) + \alpha(n)$.

    Hence, $\Phi_{q}(F) - \Phi_{q - 1}(F) \leq \alpha(n)$. The actual time of $\texttt{link}(x, y)$ is $1$ and thus the lemma follows.
\end{proof}

\begin{lemma}
    The amortised time of \texttt{find\_set} is $\O(\alpha(n))$.
\end{lemma}
\begin{proof}
    Assume the find path contains $s$ nodes :
    \begin{enumerate}
        \item There are no potential increases from \ref{lemma10}.
        \item The potential of at least $\max \left\{0, s - \alpha(n) - 2\right\}$ decreases by at least one.
              Consider all nodes $x$ on the find path with rank at least $1$ and such that there exists $y$ an ancestor of $x$ with the same \texttt{level}. There are at least $s - \left(\alpha(n) + 2\right)$ such nodes.

              Consider such a node $x$, let $k = \texttt{level}(x)$ :

              Prior to path compression, $x.p.rank \geq A_{k}^{\texttt{iter}(x)}(x.rank)$ and $y.p.rank \geq A_{k}(y.rank)$. Then, by monotony of $rank$ along a path, $y.rank \geq x.p.rank$. Let $i = \texttt{iter}(x)$. We have $y.p.rank \geq A_{k}\left(A_{k}^{i}(x.rank)\right)$

              After path compression, $x$ and $y$ have the same parent, the root so $x.p.rank = y.p.rank \geq A_{k}^{i + 1}(x.rank)$. since $x.rank$ does not change. Therefore, either \texttt{iter}($x$) or \texttt{level}($x$) increases, and hence, $\Phi_{q}(x) \leq \Phi_{q - 1}(x) - 1$ from \ref{lemma10}.
    \end{enumerate}
    As a corollary, we obtain that $\Phi_{q}(F) \leq \Phi_{q - 1}(F) - \max \left\{0, s - \left(\alpha(n) + 2\right)\right\}$. Since the actual time of \texttt{find\_path} is $s$, the amortised time is $\O\left(\alpha(n) + 2\right)$.
\end{proof}

Finally :

\begin{theorem}
    A sequence of $m$ \texttt{make\_set, union\_set, find\_set}, out of which $n$ are \texttt{make\_set} takes $\O\left(m \alpha(n)\right)$ time. In other words, one operation takes $\O(\alpha(n))$ amortised time.
\end{theorem}

\part[DFS-Matroids-MST]{Lecture 7 : 16/11}  % By Pierre Aboulpierre
\localtableofcontents
\section{Definitions}
All along the course :
\begin{definition}
    \begin{itemize}
        \item $n$ is the number of vertices
        \item $m$ is the number of edges
        \item An algorithm in $\O(n + m)$ is said to be linear
        \item $K_{n}$ is the complete graph on $n$ vertices
        \item $G = \left(U, V, E\right)$ is a bipartite graph means $U$ and $V$ is the partition.
        \item $K_{a, b}$ denotes the complete bipartite graph with parts of size $a$ and $b$.
        \item $H$ is a subgraph of $G$ if $H$ can be obtained from $G$ by deleting vertices and edges.
        \item $H$ is a induced subgraph of $G$ if $H$ can be obtained from $G$ by deleting vertices.
        \item $H$ is a subdivision of $G$ if $H$ can be obtained  from $G$ by subdividing some edges.
        \item $H$ is a minor of $G$ if $H$ can be obtained from $G$ by deleting vertices, edges and contracting edges.
        \item A sink is a vertex of outdegree 0
    \end{itemize}
\end{definition}
\begin{table*}
    \caption{Some Graph Parameters}
    \begin{tabular}{cc}
        $\delta(G)$  & minimum degree                    \\
        $\Delta(G)$  & maximum degree                    \\
        $\omega(G)$  & clique number                     \\
        $\alpha(G)$  & size of a maximum independent set \\
        $\chi(G)$    & chromatic number                  \\
        $\tau(G)$    & vertex cover                      \\
        $\kappa(G)$  & vertex connectivity               \\
        $\lambda(G)$ & edge connectivit                  \\
        $\mu(G)$     & size of a maximum matching        \\
        Girth        & length of a shortest cycle        \\
    \end{tabular}
\end{table*}


\begin{definition}[Generic Graph Search]
    The problem is to visit the vertices following the edges of the graph. A search also outputs a search tree, and an ordering on the vertices.
\end{definition}

\section{BFS, DFS and Graph Encoding}
\subsection{Properties}
\begin{definition}
    We will denote by $v.d$ the time at which a vertex is discovered, and $v.f$ the time at which a vertex is finished, i.e. all of its neighbour are discovered.
\end{definition}
\begin{theorem}[Parenthesis Theorem]
    In any DFS of a graph, for any two vertices $u$ and $v$ either :
    \begin{itemize}
        \item $\left[u.d, u.f\right]\cap \left[v.d, v.f\right] = \emptyset$ and $u$ is not a descendant nor an ancestor of $v$.
        \item $\left[u.d, u.f\right] \subset \left[v.d, v.f\right]$ and $u$ is a descendant of $v$
        \item $\left[v.d, v.f\right] \subset \left[u.d, u.f\right]$ and $v$ is a descendant of $u$
    \end{itemize}
\end{theorem}

\begin{corollary}
    A vertex $v$ is a descendant of $u$ in the DFS forest if and only if $u.d < v.d < v.f < u.f$.
\end{corollary}

\begin{theorem}[White-path Theorem]
    In a DFS forest of a digraph, a vertex $v$ is a descendant of a vertex $u$ if and only if at time $u.d$, there is a $(u, v)$-path made of undiscovered vertices.
\end{theorem}

\begin{definition}
    We define arc types in terms of the DFS forest $G_{\pi}$ :
    \begin{itemize}
        \item Tree arcs are arcs of $G_{\pi}$
        \item Forward arcs are arcs $uv$ such that $u$ in an ancestor of $v$
        \item Back arcs are arcs $uv$ such that $u$ is a descendant of $v$
        \item Cross arc are arcs $uv$ such that $u$ is not an ancestor of $v$ and $v$ is not an ancestor of $u$.
    \end{itemize}
\end{definition}

\begin{proposition}
    A digraph has a directed cycle if and only (any) DFS produces a back arc.
\end{proposition}

\subsection{Topological Sort}
\begin{definition}
    A topological ordering of the vertices of a digraph is a labeling $f$ such that $uv \in E only if f(u) < f(v)$.
\end{definition}

\begin{theorem}
    $G$ has a topological ordering if and only if it is a DAG.
\end{theorem}

\begin{theorem}
    The problem of finding a topological is solved in linear time using a DFS.
\end{theorem}

\begin{lemma}
    A digraph is acyclic if and only if a DFS of the graph yields no back edges.
\end{lemma}

\subsection{Strongly Connected Components}
\begin{definition}
    A strong connected component (scc) of a directed graph $G$ is a maximum set of vertices such that every pair of vertices as a directed path from the first to the second and from the second to the first. \\ We define $G^{SCC}$ the quotient graph for $u \mathcal{R} v \Leftrightarrow u \leadsto v \wedge v \leadsto u$.
\end{definition}

\begin{proposition}
    $G^{SCC}$ is a DAG.
\end{proposition}

\begin{lemma}
    Let $C$, $C^{'}$ be two scc of a digraph. Let $u, v \in C$ and $u^{'}, v^{'} \in C^{'}$. If $u \leadsto u^{'}$ then there is no path from $v^{'}$ to $v$.
\end{lemma}

\begin{definition}
    We extend the discovery and finishing times to sets of vertices :
    \begin{itemize}
        \item $d(S) = \min \left\{u.d \mid u \in S\right\}$
        \item $f(S) = \max \left\{u.f \mid u \in S \right\}$
    \end{itemize}
\end{definition}

\begin{lemma}
    Let $C$ and $C^{'}$ be two scc and let $u \in C$, $v \in C^{'}$ such that $uv \in E$. Then $f(C) > f(C^{'})$.
\end{lemma}

\begin{definition}
    We define $G^{R} = \left(V, E^{R}\right)$ where $E^{R} = \left\{uv \mid vu \in E\right\}$.
\end{definition}

\begin{proposition}
    $G^{R}$ has the same scc as $G$
\end{proposition}

\begin{algorithm}
    \caption{Kosaraju's Two Pass Algorithm}
    \begin{algorithmic}
        \State Call $\texttt{DFS}(G$)
            \State Compute $G^{R}$
            \State Call $\texttt{DFS}(G^{R})$ considering vertices in order of decreasing $u.f$
            \Return Output the vertices of each tree in the DFS forest computed by $\texttt{DFS}(G^{R})$
    \end{algorithmic}
\end{algorithm}

\begin{theorem}
    Kosaraju's two pass algorithm computes the scc in linear-time.
\end{theorem}


\section{Minimum Weighted Spanning Tree - MWST}
\subsection{Trees}
\begin{definition}
    A tree is a connected acyclic graph. A forest is an acyclic graph. A leaf of a tree is a vertex of degree $1$.
\end{definition}

\begin{proposition}
    \begin{itemize}
        \item Every tree contains at least two leaves
        \item Every tree $T$ satifies : $\abs{E(T)} = \abs{V(T)} - 1$.
        \item Every subgraph of a tree is a forest.
    \end{itemize}
\end{proposition}

\begin{proposition}
    Let $G = (V, E)$ be a graph. The following are equivalent :
    \begin{itemize}
        \item $G$ is a tree
        \item Any two vertices in $G$ are linked by a unique path
        \item $G$ is connected and $\abs{E} = \abs{V} - 1$
        \item $G$ is acyclic and $\abs{E} = \abs{V} - 1$
        \item $G$ is connected but for every edge $uv$, $G - \{uv\}$ has exactly two connected components, one containg $u$ and the other $v$.
        \item $G$ is acyclic but if any edge is added to $G$, the resulting graph contains a unique cycle going through this added edge.
    \end{itemize}
\end{proposition}

\begin{definition}
    $\omega : E \to \R$ is called a weight function. A spanning tree is a subgraph of $G$ that is a tree and contains all vertices of $G$. The weight of a subgraph is the sum of the weights of its edges.
\end{definition}


\begin{proposition}[Cut and Paste Technic]
    Let $G = (V, E)$ and let $T$ be a spanning tree of $G$. Let $uv \in E(G) - T$ and let $T_{uv}$ be the unique path linking $u$ and $v$ in $T$. Then for every edge $xy$ of $T_{uv}$ $T\setminus \{xy\} \cup \{uv\}$ is a spanning tree of $T$.
\end{proposition}
\begin{proof}
    Since $T_{uv}$ is in the unique path of $T$ linking $u$ and $v$, removing any edge $xy \in T_{uv}$ from $T$ breaks $T$ into two connected components, one containing $u$ and the other $v$. Adding $uv$ reconnects the two parts and form a new spanning tree $T \setminus \{xy\} \cup \{uv\}$.
\end{proof}

\begin{theorem}
    Let $G = \left(V, E, \omega\right)$. A spanning tree $T$ of $G$ is a minimum spanning tree if and only if for every edge $e \in E \setminus T$ :
    \[
        \omega(e) \geq \omega(f) \forall f \in \text{ the unique cycle of } T \cup \{e\}
    \]
\end{theorem}

\begin{definition}
    A set of edges $A$ is said to be promising if it can be completed into a minimum spanning tree. At each step, we determine an edge $uv$ that we can add to $A$ without violating this invariant, that is $A \cup \{uv\}$ is still a promising set.
    Such an edge $uv$ is said to be safe with respect to $A$.
\end{definition}

The algorithm works as follows : We start with an empty set, and while our set is not a spanning tree, we find a safe edge and add it to the set.
How to find safe edges ?

\begin{definition}
    A cut $\left(S, V - S\right)$ is a partition of $V$. An edge of the cut or crossing edge is an edge with one end in $S$ and the other in $V - S$. A cut respects a set of edges if no edge of the set is a crossing edge of the cut.
\end{definition}

\begin{proposition}
    Let $A$ be a promising set of edges. Let $\left(S, V - S\right)$ be a cut respecting $A$ and $uv$ be a lightest crossing edge. Then $uv$ is safe.
\end{proposition}
\begin{proof}
    Let $T$ be a MST containing $A$. If $T$ contains $uv$ we are done, so assume it does not. Let $T_{uv}$ be the unique path linking $u$ and $v$ in $T$. Since $uv$ is a crossing edge of $(S, V - S)$, $T_{uv}$ contains an edge $xy$ of $(S, V - S)$. By the cut and paste technique, $T^{'} = (T \setminus \{xy\} ) \cup \{uv\}$ is a spanning tree of $G$. By hypothesis, $uv$ is a lightest edge of $(S, V - S)$, so $\omega(uv) \leq \omega(xy)$ and thus: \[\omega(T^{'}) = \omega(T) + \omega(uv) - \omega(xy) \leq \omega(T)\]
    So $T^{'}$ is a MST containing $uv$
\end{proof}

\subsection{KRUSKAAAAAAAAL}
Kruskal algorithm grows a promising forest $A$ :
\begin{itemize}
    \item Sort the edges by non-decreasing order of weight
    \item Start with the empty set
    \item Add a minimum weighted edge that does not create a cycle which is equivalent to add a minimum weighted edge that connects two connected components of the growing forest.
\end{itemize}

\begin{algorithm}
    \caption{Kruskal}
    \begin{algorithmic}
        \Input \ \ $G = \left(V, E, \omega\right)$
        \EndInput
        \State $A \gets \emptyset$
        \For {$u \in V$}
        \State $\texttt{Make-Set}(u)$
        \EndFor
        \State $\texttt{Sort}(E)$ \Comment{Sort by non-decreasing order of weight}
        \For {$uv \in E$}
        \If {$\texttt{Find}(u) \neq \texttt{Find}(v)$}
        \State $A \gets A \cup \{uv\}$
        \State $\texttt{Union}(u, v)$
        \EndIf
        \EndFor
        \Return \ \ A
    \end{algorithmic}
\end{algorithm}

\begin{theorem}
    The Kruskal algorithm answers the problem in $\O(m \log(n))$ time
\end{theorem}

\begin{proof}
    \begin{itemize}
        \item Complexity : from union-find, we do $\O(m)$ \texttt{Find-Set} and $\O(n)$ \texttt{Union} and \texttt{Make-Set} in $\O(m + n)\alpha(n)^{4}$. As we do one sort, the total complexity is $\O(m\log(n))$.
        \item Correctness : Comes from the safe property, by considering the loop invariant : Prior to each iteration, $A$ is promising.
    \end{itemize}
\end{proof}

\subsection{Prim}
This algorithm works like Kruskal's, but instead of growing a forest, we grow a tree $A$. At each iteration, we add the lightest edge with exactly one extremity in $V(A)$.
\begin{definition}
    A Priority Queue is a data structure that maintains a set $S$ of elements, each with an associated value called a key. TIt supports the following :
    \begin{itemize}
        \item $\texttt{Insert}(S, x)$
        \item $\texttt{Minimum}(S)$
        \item $\texttt{Extract-Min}(S)$
        \item $\texttt{Decrease-Key}(S, x, k)$ does $x.key \gets k$.
    \end{itemize}
\end{definition}
\begin{remark}
    This is for a min-queue, we can also do a max-queue.
\end{remark}

We implement priority queues using a heap.
\begin{definition}
    A max-heap is a complete binary tree where the key of a node is larger than the keys of its children. It can be implemented in an array : $\texttt{Parent}(i) = \lfloor \frac{i}{2}\rfloor, \texttt{Left}(i) = 2 i + 1, \texttt{Right}(i) = 2i + 2$. It supports the same operations as a max-queue.
\end{definition}

\begin{theorem}
    Using a binary heap, insertion is done in $\log N$, max-retrieval is done in constant time and max-deletion is done in $\log N$.
\end{theorem}
\begin{proof}
    See implementation.
\end{proof}

\begin{algorithm}
    \caption{PRIM}
    \begin{algorithmic}
        \Input \ $G = \left(V, E, \omega\right), r$
        \EndInput
        \For {$v \in V$}
        \State $v.\pi \gets \texttt{NIL}$ and $v.key \gets + \infty$
        \EndFor
        \State $r.key \gets 0$
        \State $Q \gets V$
        \While {$Q \neq \emptyset$}
        \State $u \gets \texttt{Extract-Min}(Q)$
        \For $v \in \texttt{Adj}[u]$
        \If $v \in Q \wedge v.key > \omega(uv)$
        \State $v.key \gets \omega(uv)$
        \State $v.\pi \gets u$
        \EndIf
        \EndFor
        \EndWhile
    \end{algorithmic}
\end{algorithm}

\begin{theorem}
    This algorithm solves the problem in $\O(m + n\log n)$.
\end{theorem}

\begin{proof}
    \begin{itemize}
        \item Correctness comes from the safe property
        \item Complexity is in $\O(m\log n)$ with a min-heap and $\O(m + n\log n)$ with a Fibonacci Heap
    \end{itemize}
\end{proof}

\section{Matroids and Greedy Algorithms}
\subsection{Definitions}
\begin{definition}
    Let $E$ be a set of elements and $\mathcal{I} \subseteq 2^{E}$. $\left(E, \mathcal{I}\right)$ is a hereditary set system if it satisfies
    \begin{itemize}
        \item M1 : $\emptyset \in \mathcal{I}$
        \item M2 : If $X \subseteq Y \in \mathcal{I}$, $X \in \mathcal{I}$.
    \end{itemize}
    \begin{itemize}
        \item $E$ is called the ground set.
        \item Sets in $\mathcal{I}$ are called independent sets.
        \item Maximal independent sets are called bases.
        \item For $X \subseteq E$, the rank of $X$ denoted by $rk(X)$ is the size of a maximum independent set included in $X$.
        \item Sets not in $\mathcal{I}$ are called dependent sets.
        \item Minimal dependent sets are called circuit.
    \end{itemize}
\end{definition}

Given a positive weight function, we want to find a maximum independent set such that its weight is maximum. For example :

\begin{itemize}
    \item MST : $\mathcal{I} = \left\{I \subseteq E \mid I \text{ is a forest } \right\}$
    \item TSP, SPP, Maximum Matching Problem, Maximum Stable Set, Knapsack are also of the sort.
\end{itemize}


\begin{theorem}[Rado - Edmonds]
    The greedy algorithm for this problem is optimal for any weight function if and only if $\left(E, \mathcal{I}\right)$ is a matroid.
\end{theorem}

\begin{definition}
    A matroid is a hereditary set system $E, \mathcal{I}$ such that : if $X, Y \in \mathcal{I}$ and $\abs{X} > \abs{Y}$, there is $x \in X \setminus Y$ such that $Y \cup \{x\} \in \mathcal{I}$
\end{definition}

\part[Parametrized Complexity]{Lecture 9 : 21/12}
\section{Complexity}
\subsection{Definitions}
We usually aim for poly-time algorithm in worst case. Since many problems of interest are $NP$-hard, we want to find better complexity than $\O(c^{n})$ for them, as poly-time in unlikely. \\
For $NP$-hard problems, there is no algorithm that solves all instances optimally in poly-time. We will now look to measure the complexity in term of the input size and a parameter of the input.

For example, \textsc{Vertex Cover} is $NP$-complete, and looks for a graph $G$ and an integer $k$, is it possible to cover the edges with $k$ vertices ? In Brute-Force, there are $\O(n^{k})$ possibilities but a $\O(2^{k}n^{2})$ algorithm exists. Yet, for $k$-\textsc{Independent Set}, there are no $n^{o(k)}$ algorithm known.

\begin{definition}
    A parametrized algorithmic problem is a problem where a certain parameter $k$ is given in addition to the input. The complexity is studied as a function of $n$ and $k$. $k$ can be the size of the solution, or an implicit parameter of the input graph.
\end{definition}

We then have three main possibilities :
\begin{enumerate}
    \item Either the problem is already hard for fixed $k$ : $\chi(G) \leq k$.
    \item Or the problem is $NP$-hard for $k$ in the input but polynomial for fixed $k$ : $\alpha(G) \leq k$. We say the problem is $XP$
    \item Or it is FPT for $k$, there is an algorithm in $\O(f(k)\times n^{\O(1)})$
\end{enumerate}

Formally :
\begin{definition}
    Let $\Sigma$ be a finite alphabet used to encode the input. A parametrized problem is a set $\mathcal{P} \subseteq \Sigma^{\star} \times \N$. The set $\mathcal{P}$ contains $(x, k)$ if the answer to the question encoded by $x, k$ is \texttt{true}. $k$ is the parameter.
\end{definition}

\begin{definition}
    \begin{itemize}
        \item A parametrized problem is Fixed-Parameter Tractable if there is an algorithm deciding $\mathcal{P}$ in time $f(k)n^{c}$ where $f$ is computable and $c$ is constant.
        \item A parametrize problem $\mathcal{P}$ is $XP$ if there is an algorithm deciding $\mathcal{P}$ in time $n^{f(k)}$ for some computable $f$ and constanc $c$.
    \end{itemize}
\end{definition}

The $W$ hierarchy is a hierarchy inside of $XP$ that looks at the parametrized problems in the same way as $P$ hierarchy.

\begin{openpb}
    $n$-variable $3$-SAT cannot be solved in time $2^{o(n)}$
\end{openpb}

\subsection{Clique by Degree}
We define :
\begin{definition}
    Given a graph $G$ with maximum degree $\Delta$ and an integer $k$, does $G$ have a clique of size at least $k$ ?
\end{definition}

This is FPT since looking at the neighbourhoods of each vertex gives us $\O(2^{\Delta}n)$ time complexity. There might be no FPT solution for \textsc{Clique} parametrized by $k$.

\section{Algorithmic Techniques}
\subsection{Branching Method}

\subsubsection{$k$-\textsc{Vertex Cover}.}

\begin{definition}
    The main idea of the branching method is to reduce the problem to solving a bounded number of problems with parameter $k' < k$. Since the parameter decreases in every recursive call, the depth of the search tree is at most $k$. The size of the search tree branching into $c$ directions in $c^{k}$.\\
    We denote $\O(\alpha^{k}n^{\O(1)})$ by $\O^{\star}(\alpha^{k})$
\end{definition}

Then, for \textsc{Vertex Cover} we branch on vertices of degree at least $3$ :
\begin{itemize}
    \item For each vertex $u$ of degree at least $3$, either $u$ is in the solution, or all the neighbours of $u$ are in the solution.
    \item If every vertex has degree at most $2$, we can solve \textsc{Vertex Cover} in poly-time because the graph is the disjoint union of paths and cycles. These graphs will be the leaves of our search-tree.
\end{itemize}


\begin{proposition}
    The above algorithm gives us a solution to \textsc{Vertex Cover} in $\O^{\star}(1.4656^{k})$
\end{proposition}
\begin{proof}
    Let $T(k)$ be the number of leaves in the search tree and $T(k) = 0$ if $k \leq 1$. \\
    We have $T(k) \leq T(k - 1) + T(k - 3)$. We want to find $c \geq 1$ such that $T(k) \leq c^{k}$ for all $k$. \\
    We see that $c$ such that $c^{k} \geq c^{k - 1} + c^{k - 3}$. In particular, $c^{3} - c^{2} - 1 \geq 0$. Then $c = 1.4656$ is valid.
\end{proof}
\begin{remark}
    There exists an algorithm in $\O^{\star}(1.2738^{k})$.
\end{remark}

\begin{definition}
    We define the branching vector of an algorithm as the vector with the decrease parameters.
\end{definition}
Here it was $(1, 3)$. If we find an algorithm that branches for $(2, 5, 6, 6, 7, 7)$, we will get $\O^{\star}(1.4483^{k})$

\subsubsection{$k$-\textsc{Feedback Vertex Set}}
\begin{definition}
    Given a multigraph $G$ and an integer $k$, find a set $S$ of at most $k$ vertices such that $G \setminus S$ is a forest.
\end{definition}
A feedback vertex set is a set of vertices that hits every cycle of the graph.

Here, at least one vertex of each cycle must be in the solution, but the size of a cycle can be arbitrarily large. We are going to identify a set of $\O(k)$ vertices such that any size $k$-FVS has to contain one of these vertices and branch on it. To do so, we need reduction rules.

\begin{proposition}[Reduction Rules for FVS]
    \begin{enumerate}
        \item If there is a loop at $v$, then delete $v$ and decrease $k$ by one.
        \item If there is an edge of multiplicity large than $2$, reduce its multiplicity to $2$
        \item If there is a vertex $v$ of degree $\leq 1$, delete it.
        \item If there is a vertex $v$ of degree $2$, delete $v$ and add an edge between the neighbours of $v$.
    \end{enumerate}
    Applying the reduction rules returns a graph $G^{'}, k'$ with no loops, edge multiplicty > 2, minimum degree $3$.

    We have $G, k$ a YES instance if and only if $G^{'}, k'$ is a YES instance.
\end{proposition}

\begin{lemma}
    Let $G$ be a graph with minimum degree $3$, let $V_{3k}$ be the $3k$ largest degree vertices. Then every $FVS$ of size at least $k$ contains at least a vertex of $V_{3k}$.
\end{lemma}
\begin{proof}
    Suppose there is $S$ a solution disjoint frome $V_{3k}$ and let $d = \min d(V_{3k})$. Let $X = V(G) \setminus \left(S \cup V_{3k}\right)$.
    \[
        \sum_{v \in X \sqcup V_{3k}} d(v) \geq 3 \abs{X} + 3kd
    \]
    But $G[X \cup V_{3k}]$ is a forest so its number of edges is $\leq \abs{X} + 3k - 1$. So the sum of degrees in it is at most $2\abs{X} + 6k - 2$. \\
    Thus we get the number of edges to be :
    \begin{itemize}
        \item $\geq 3kd + 3\abs{X} - \left(2 \abs{X} + -k - 2\right) > 3kd - 6k$
        \item $\leq dk$ because $S$ has $k$ vertices of degree at most $d$.
    \end{itemize}
    So we would have $2kd - 6k < 0$ which is impossible since $d \geq 3$.
\end{proof}

We then get an algorithm to compute a solution by applying the reduction rules and branching (in polynomial time) on each vertex in $V_{3k}$. Since we branch into $3k$ directions, we get a time complexity in $\O^{\star}\left((3k)^{k}\right)$

\subsection{Kernelization}
\subsubsection{Definitions}
Kernelization is a method for parametrized preprocessing. We want to reduce the size of the instance $(x, k)$ to an equivalent instance with size bounded by $f(k)$.
\begin{definition}
    Let $\mathcal{P}$ be a parametrized problem and $f$ a computable function.\\
    A kernel of size $f(k)$ is an algorithm that, given $(x, k)$, runs in polynomial-time in $\abs{x} + k$ and outputs an instance $x', k'$ such that :
    \begin{itemize}
        \item $x, k \in \mathcal{P} \Leftrightarrow x', k' \in \mathcal{P}$
        \item $\abs{x^{'}} \leq f(k)$ and $k' \leq k$.
    \end{itemize}
    We say the kernel is polynomial is $f$ is polynomial.
\end{definition}

\begin{theorem}
    A parametrized problem is FPT if and only if it is decidable and has a kernel.
\end{theorem}
\begin{proof}
    \begin{itemize}
        \item If the problem has a kernel, we reduce the size of the instance in poly-time and bruteforce on it.
        \item If the problem is solved in $f(k)\cdot \abs{x}^{c}$.\\
        If $\abs{x} \leq f(k)$ then we have our kernel. Else, if $\abs{x} \geq f(k)$ then we can solve the problem in time $f(k) \cdot \abs{x}^{c} \leq \abs{x}^{c + 1}$
    \end{itemize}
\end{proof}

The real issue is to know which problems have polynomial kernels. 

\subsubsection{\textsc{Vertex Cover}}
We define two reduction rules for $k$-\textsc{Vertex Cover} : 
\begin{proposition}
    \begin{itemize}
        \item If $v$ has degree $0$, then reduce to $(G - v, k)$
        \item If $v$ has degree $k + 1$, then reduce to $(G - v, k - 1)$.
    \end{itemize}
\end{proposition}

\begin{lemma}
    If $(G, k)$ is a YES instance for $k$-\textsc{Vertex Cover} on which reduction rules cannot be applied, $G$ has at most $k^{2}$ edges and $k^{2} + k$ vertices.
\end{lemma}
\begin{proof}
    Let $S$ be a vertex cover of $G$ of size at most $k$. Each vertex hits at most $k$ edges because of the second rule. So there is at most $k^{2}$ edges. Each vertex is either in $S$ or is one of the at most $k$ neighbours of a vertex in $S$ so there are at most $k^{2} + k$ vertices. 
\end{proof}

\begin{proposition}[Kernelization for \textsc{Vertex Cover}]
    \begin{itemize}
        \item Apply reduction rules exhaustively. 
        \item We get an equivalent instance $(G^{'}, k')$ with $k' \leq k$.
        \item If $\abs{E(g)} > k'^{2}$ $\abs{V}(G) > k'^{2} + k'$ return NO
        \item Otherwise we have a kernel of size $\O(2k^{2} + k)$.
    \end{itemize}
\end{proposition}

\begin{theorem}
    \textsc{Vertex Cover} has a kernel with at most $2k$ vertices. 
\end{theorem}
This is done using Linear Programming. 

\subsection{Color Coding}
\subsubsection{$k$-\textsc{Path}}
\begin{definition}
    Given $G, k$, decide if $G$ contains a path on $k$ vertices. 
\end{definition}

We will use a randomize\footnote{BOOOOOOOOOOOOOOOOOOOOOOOOOOOOOOOOO} algorithm that can be determinized.

We transform the problem into the following : 
\begin{itemize}
    \item Assume the vertices are colored randomly with $\left\{1, \ldots, k\right\}$. 
    \item Find a path colored $1 - 2 - \cdots - k$. 
\end{itemize}

The algorithm is the following : 
\begin{itemize}
    \item Assign color from $\left[k\right]$ to the vertices uniformly and independently at random
    \item Output YES if there is a path colored $1 - 2 - \cdots - k$ else NO
\end{itemize}
\begin{proposition}
    If $G$ has a $k$-path, the probability that this $k$-path is colored $1 - 2 - \cdots - k$ is $1/k^{k}$. So if $G$ is a YES instance the algo outputs YES with probability at least $1/(k^k)$, and if it is a NO instance, the algorithm outputs NO. Then, after $k$ repetitions, error probability is at most $1/e$. 
\end{proposition}

To find a colored path, we delete edges in the graph from vertices colored $i$ to vertices colored $j$ for $j \neq i-1, i + 1$ and orient edges towards the larger class. We can find a path in linear time with BFS. 

\begin{proposition}
    The final complexity of this algorithm is $\O(c\cdot k^{k} \cdot (n + m))$ and it has probability of success $1/e^{c}$.
\end{proposition}

We can still improve color coding by solving the existence of a colorful path instead of the existence of an ordered coloured path. Indeed, if there is $k$-path, the probability that is it colorful is $> e^{- k}$. We now only need to solve the problem $e^{k}$ times instead of $k^{k}$.\\
We can solve this problem using Dynamic Programming : we solve for each vertex $v$ and set of color $C \subseteq \onen{k}$, the value of $D(v, C)$ that is YES is there is path ending at $v$ using each colour of $C$. \\
If we denote by $\chi : V \rightarrow [k]$ the random coloring, $D(v, C)$ is YES if and only if $\chi(v) \in C$ and there is an edge $uv$ for which $D(u, C\setminus \chi(v))$ is YES. We thus can solve this in time $2^{k}\times \abs{E}$.

Finally we get the following algorithm : 
\begin{proposition}
    Repeat $e^{k}$ times : 
    \begin{enumerate}
        \item Sample a coloring $c$
        \item Check if $G$ contains a colorful $k$-path in time $\O(2^{k}\cdot \abs{E})$ and return YES if it does. 
    \end{enumerate}
    If no colorful $k$-path was found, return NO. 
\end{proposition}

\subsubsection{Derandomization}
\begin{definition}
    A family $\mathcal{H}$ of function $\onen{n} \rightarrow \onen{k}$ is a $k$-perfect family of hash functions if for every $S \subseteq \onen{n}$ with $\abs{S} = k$, there is an $h \in \mathcal{H}$ such that $h(x) \neq h(y)$ for any $x\neq y \in S$.
\end{definition}

\begin{theorem}
    There is a $k$-perfect family of functions $\onen{n} \rightarrow \onen{k}$ with size $2^{\O(k)}\log n$ and can be constructed in poly-time in the size of the family. 
\end{theorem}

Then, instead of trying $\O(e^{k})$ random colorings, we go through a $k$-perfect family $\mathcal{H}$ of functions $\onen{n} \rightarrow \onen{k}$. Then, $k$-\textsc{Path} can be solved in deterministic time $2^{\O(k)}\cdot n^{\O(1)}$.


\part[Approximation Algorithm]{Lecture 10 : 11/01}
\section{Definitions}
Suppose we want to solve an optimization problem $\mathcal{P}$. Say that a solution has a certain cost we want to minimize. Denote by $OPT(P)$ the cost of an optimal solution of $\mathcal{P}$. Suppose we have an algorithm that outputs a non-optimal solution of cost $ALG(P)$. $ALG$ has approximation ratio $\rho(n)$ if for any instance $P$ of size $n$ we have : \[ ALG(P) \leq \rho(n)\cdot OPT(P)\] 

\begin{definition}
    An algorithm has an approximation ratio of $\rho(n)$ if for any input of size $n$ the cost ALG of the solution produced by the algorithm is with a factor of $\rho(n)$ of the cost $OPT$ of an optimal solution : 
    \[
        1 \leq \max \left(\frac{OPT}{ALG}, \frac{ALG}{OPT}\right) \leq \rho(n)
    \]
\end{definition}

Many NP-Hard problems  have small constant approximation ratio. Others have approximation ratio that grows with $n$. When the best approximation ration grows linearly, there is trouble. 

An approximation scheme is an algorithm scheme is an algorithm that outputs a $(1 + \epsilon)$-approximation where $\epsilon$ is part of the input.
\begin{definition}
    An algorithm is a Polynomial Time Approximation Scheme if it satifies
    \begin{itemize}
        \item Input : An instance $P$ and $\epsilon > 0$
        \item Output : A $(1+\epsilon)$-approximation
        \item Polynomial in $n$ for every fixed $\epsilon$
    \end{itemize}
    We say it is a Fully Polynomial Time Approximation Scheme if it is polynomial in $n$ and $1/\epsilon$.
\end{definition}

\subsection{\textsc{Min Vertex-Cover}}
\begin{algorithm}
    \caption{Naïve Algorithm}
    \begin{algorithmic}
        \State For each vertex $v$
        \If {$d(v) \geq 1$}
            \State Put $v$ in the solution
            \State Delete $v$
        \EndIf
    \end{algorithmic}
\end{algorithm}
Yet, for a star graph, the algorithm returns a set of size $n-1$ while there is a solution of size $1$. 

\begin{algorithm}
    \caption{Greedy Algorithm}
    \begin{algorithmic}
        \While {there are edges}
        \State Choose Vertex with Maximum Degree
        \State Delete It
        \EndWhile
    \end{algorithmic}
\end{algorithm}

We need to find a lower bound of an optimal solution : \\
Assume there is a certain quantity $Q$ such that $Q \leq OPT$. If we can prove that $ALG \leq c \times Q$ then $ALG$ is a $c$-approximation. 
Here, a Vertex Cover is at least as large as a maximal matching. We get : 

\begin{algorithm}
    \caption{Based on Matching}
    \begin{algorithmic}
        \State $M \gets$ Maximal Matching
        \While {there are edges}
        \State Add $u$ and $v$ to the solution
        \State Delete $u$ and $v$
        \EndWhile
    \end{algorithmic}
\end{algorithm}

Here, we have a $2$-approximation algorithm. We need to see if we can improve the approximation ratio. It is tight from looking at $K_{n, n}$. \\
State of the art, the best approx algorithm known has approximation factor $\left(2 - \Theta\left(\frac{1}{\sqrt{\log n}}\right)\right)$. If Unique Game Conjecture is true, we cannot find better that $2$ for a constant factor. 

\subsection{\textsc{Travelling-Salesman Problem}}
\begin{definition}
    TSP asks the question of given a list of cities and distances between each pair of cities, find a hamiltonian cycle (that goes every vertex exactly once) of minimum cost ?
\end{definition}

We can assume WLOG that the graph is complete. We denote by $c(A)$ the weight of the path $A$. \\
\begin{theorem}
    For every constant $\rho \geq 0$, there is no polynomial-time algorithm for TSP with approximation ratio $\rho$, unless $P = NP$
\end{theorem}
\begin{proof}
    Suppose there is an algorithm $ALG$ that computes a $\rho$-approximation of TSP in polynomial-time on every input $G$.\\
    Then, let $G$ be an instance for \textsc{Hamiltonian Path}. We transform $G$ intro an instance for $TSP$ by taking a complete graph on $V(G)$ with cost \[c(uv) = \begin{cases}
        1 & \text{ if } uv\in E(G)\\
        \rho\abs{V} &\text{ otherwise}
    \end{cases}\]
    Let $OPT$ be the cost of an optimal tour of $G^{'}, c$. Observe that, if $G$ has an hamiltonian cycle, then $OPT = \abs{V(G)}$, else $OPT > \rho \times \abs{V(G)}$.
\end{proof}

\begin{definition}
    \textsc{Metric Travelling-Salesman Problem} is \textsc{TSP} where the weight function $c$ of the graph is a metric. 
\end{definition}

\begin{theorem}
    \textsc{MTSP} has a $2$-approximation algorithm. 
\end{theorem}
\begin{proof}
    We have $c{TSP} > c{MST}$. Indeed, removing an edge $TSP$ gives a spanning tree.
    \begin{algorithm}
        \caption{Metric TSP using MSTs}
        \begin{algorithmic}
            \State $M \gets MST$ (using Prim or Kruskal)
            \State\Return Preorder Walk of $M$
        \end{algorithmic}
    \end{algorithm}
    This is a two approx : If $T$ is a 
\end{proof}

\subsection{\textsc{Set-Covering}}
\begin{definition}
    Given a hypergraph $\left(V, \mathcal{F}, c\right)$ where $V$ is finite a set of elements, $\mathcal{F}$ is a set of subsets of $V$ and $c$ is a cost function, output a cost minimum subset of $\mathcal{F}$ that covers all elements of $V$. 
\end{definition}
First, a greedy algorithm. Here, the criteria is basing on the number of uncovered vertices.\\
\begin{algorithm}
    \caption{Greedy Algorithm}
    \begin{algorithmic}
        \State $U \gets V$
        \State $\mathcal{C} \gets \emptyset$
        \While {$U \neq \emptyset$}
            \State Select $S \in \mathcal{F}$ that maximizes $\abs{S \cap U}$
            \State $\mathcal{C} \gets \mathcal{C} \cup \{S\}$
            \State $U \gets U \setminus S$
        \EndWhile
    \end{algorithmic}
\end{algorithm}
Yet, if there is a set in $\mathcal{F}$ of weight $10^{10}$ that is covered by two sets of weight $1$, the algorithm will choose the first option. So, we modify the greedy criteria to act on the average cost of newly covered vertices : $c(S) /\abs{S \cap U}$
\begin{algorithm}
    \caption{\textsc{Greedy Set-Cover}}
    \begin{algorithmic}
        \State $U \gets V$
        \State $\mathcal{C} \gets \emptyset$
        \While {$U \neq \emptyset$}
            \State Select $S \in \mathcal{F}$ that minimizes $\frac{c(S)}{\abs{S \cap U}}$
            \State $\mathcal{C} \gets \mathcal{C} \cup \{S\}$
            \State $U \gets U \setminus S$
        \EndWhile
    \end{algorithmic}
\end{algorithm}
\begin{theorem}
    \textsc{Greedy-Set-Cover} is a $\log(n)$-approximation.
\end{theorem}
\begin{proof}
    Define $price(v) = c(S)/\abs{S \cap U}$ where $S$ is the hyperedge that covers $v$ for the first time. \\
    The total cost of a solution is $\sum_{v \in V} price(v)$. Number the elements of $V$ by order of choose, with $v_{n}$ being the first chose one. We have : 
    \begin{lemma}
        For each $k \in \onen{n}, price(v_{k}) \leq OPT/k$
    \end{lemma}
    \begin{proof}
        Assume the algorithm is going to choose a set $S$ that is going to cover $v_{k}$ for the first time. Let $U$ be the set of uncovered elements at this moment. So $\{v_{k}, \ldots, v_{1}\} = U$.
        Since an optimal solution covers $U$ with at most $OPT$, there must be a set $S$ with cost-effectiveness at most $OPT/\abs{U}$. Hence, by the greedy price criteria
        \[
            price(v_{k}) \leq \frac{OPT}{\abs{U}} \leq \frac{OPT}{k}
        \]
    \end{proof}
    Hence, the total cost is : 
    \[
        \sum_{k = 1}^{n}price(v_{k}) \leq OPT \sum_{k = 1}^{n} \frac{1}{k} \leq OPT\cdot (\log n + 1)
    \]

    
\end{proof}
\begin{theorem}
    \textsc{Greedy-Set-Cover} does not have an approximation ration of $\left(1 - o(1)\right)\ln n$. 
\end{theorem}

\subsection{Layering and \textsc{Weighted Vertex-Cover}}
\begin{definition}
    Given a graph and a weight function, find a minimum weight vertex cover.
\end{definition}

\begin{definition}
    A degree weighted function of a graph $G$ is a function $\omega : V \rightarrow \N$ such that 
    \[
        \exists c, \forall v, \omega(v) \leq c \cdot \deg v
    \]
\end{definition}

\begin{lemma}
    Let $G$ be a graph with weight function $\omega$. If $\omega$ is a degree weighted function, then $\omega(V) \leq 2\cdot OPT$
\end{lemma}
\begin{proof}
    We have : 
    \[
        \omega(V) =  \sum_{v \in V} \omega(v) \leq \sum_{v \in V}c\deg(v) \leq 2c\abs{E}
    \]
    If $S^{\star}$ is an optimal solution, 
    \[
        OPT = \omega(S^{\star}) = \sum_{v \in S^{\star}} \omega(v) = c \cdot \sum_{v \in S^{\star}} \deg v
    \]
    Since $S^{\star}$ covers all edges we have \[\sum_{v \in S^{\star}} \deg v \geq \abs{E}\]
\end{proof}

\begin{algorithm}
    \caption{Layering Algorithm$(G, \omega)$}
    \begin{algorithmic}
        \While {$G \neq \emptyset$}
            \State Remove $D = \left\{v \mid \deg v = 0\right\}$
            \State Compute $c = \min \left\{\omega/\deg(v)\right\}$
            \State Compute degree-weighted function $t = c \cdot \deg(v)$
            \State Compute the residual weight function $\omega_{1} = \omega - t = \omega- c \cdot \deg$ 
            \State Pick the set $W = \left\{v \mid \omega_{1}(v) = 0\right\}$
            \State $G_{1} = G - \left(D \cup W\right)$
            \State $G \gets G_{1}, \omega \gets \omega_{1}$.
        \EndWhile
    \end{algorithmic}
\end{algorithm}

\begin{theorem}
    The layering algorithm is a $2$-approximation
\end{theorem}
\begin{proof}
    It is clearly correct.\\
    Let $S = W_{0} \cup \ldots \cup W_{k}$ be the solution output by the algorithm. Observe that $\omega_{l}(v) = \omega_{0}(v) - \sum_{i = 0}^{l- 1}t_{i}(v)$. Let $v \in S$. If $v \in W_{l}$, $w_{0}(v) = \sum_{i = 0}^{l - 1}t_{i}(v)$.\\
    Let $S^{\star}$ be an optimal vertex cover of $G$. For each layer $i$, $S^{\star} \cap V\left(G_{i}\right)$ is a vertex cover of $G_{i}$. Therefore by the degree weighted function Lemma : 
    \[
        t_{i}\left(S\cap V(G_{i})\right) \leq 2 t_{i}\left(S^{\star} \cap V(G_{i})\right)
    \]
    Hence : 
    \[
        \omega(S) = \sum_{i = 0}^{k} t_{i}\left(S \cap V(G_{i})\right) \leq 2 \sum_{i = 0}^{k} t_{i}(S^{\star} \cap V(G_{i})) \leq 2 \omega(S^{\star})
    \]
\end{proof}
\part[Linear Programming]{Lecture 11 : 18/01}
\section{Definition}
\subsection{Linear Programming}
\begin{definition}
    A system of linear equation is of the form : 
    \[
        \begin{aligned}
            a_{11} x_{1} + \ldots + a_{1n} x_{n} &= b_{1}\\
            \vdots + \ldots + \vdots &= \vdots \\
            a_{n1}x_{1} + \ldots + a_{nn}x_{n} &= b_{n}
        \end{aligned}
    \]
    We write it as $Ax = b$ where $A$ is a matrix. 
\end{definition}
\begin{definition}
    A linear program is made of $n$ decision variables $x_{1}, \ldots, x_{n} \in \R$, $m$ linear constraints
    \[
        \sum_{j = 1}^{n} a_{ij} \star b_{i}
    \]
    where $\star \in \{\leq, \geq, =\}$; and a function we want to maximise. 
\end{definition}

\begin{proposition}
    Linear Programming and Linear Algebra work in a similar way : 
    \begin{center}
        \begin{tabular}{p{.15\textwidth}p{.25\textwidth}p{.25\textwidth}c}
            & Basic Problem & Algorithm & Solution Set\\
            \midrule
            Linear \newline Algebra & System of \newline Linear Equations & Gaussian Elimination & Affine Subspace\\
            \midrule
            Linear \newline Programming & System of \newline Linear inequalities & Simplex Method & Convex Polyhedron
        \end{tabular}
    \end{center}
\end{proposition}

Linear Programming allows to model most of the problems we have seen in this class. Despite this generality, Linear Programs can be solved efficiently, both in theory (weak polynomial time) and in practice.
\begin{definition}
    If we allow the variable to be only in $\N$ (instead of $\R$), we call it Integer Linear Programming and it becomes NP-hard.
\end{definition}

\subsection{First Example}
Suppose a factory makes tables and chairs. Each table returns a profit of $200\$$ and each chair a profit of $\$100$. Each table takes $1$ unit of metal and $3$ units of wood and each chair takes $2$ units of metal and $1$ unit of wood. The factory has $6k$ units of metal and $9k$ units of wood. We want to maximize profit.\\
We can model this problem using linear Programming, let $x_{1}$ be the number of tables and $x_{2}$ be the number of chairs 
\[
    \begin{aligned}
        3x_{1} + x_{2} &\leq 9\\
        x_{1} + 2x_{2} &\leq 6\\
        \textnormal{Maximize } 2x_{1} + x_{2}
    \end{aligned}    
\]

\subsection{Flow}
\begin{proposition}
    We are given $G = (V, E, s, t, c)$ with: 
    \begin{itemize}
        \item A directed Graph $G = V, E$
        \item A source vertex $s \in V$
        \item A tap vertex $t\in V$
        \item A non-negative integral capacity function $c$. 
    \end{itemize}
    Introduce one decision variable for each edge $\{f_{e}\}_{e \in E(G)}$
    \[
        \begin{aligned}
            f_{e} \leq 0
        \end{aligned}
    \]
\end{proposition}

\section{Integer Linear Programming and Relaxation}
\subsection{Definition}
\begin{definition}
    The relaxation of an integer linear program is the problem that arises by removing the integrality constraint of each variable.
\end{definition}

\subsection{\textsc{Vertex Cover}}
\begin{proposition}
    Introduce $x_{v} \in \{0, 1\}$ for each $v \in V(G)$. Setting $x_{v} = 0$ means that $x_{v}$ is not in the solution.\\
    We want to minimize $\sum_{v \in V}x_{v}$ subject to 
    \[
        \begin{aligned}
            x_{u} + x_{v} &\geq 1 \text{  for all } uv \in E(G)\\
            x_{v} &\in \{0, 1\}  \text{  for all } v \in V(G)
        \end{aligned}
    \]
    Its relaxation is called Fractional Vertex Cover: 
    We want to minimize $\sum_{v \in V}x_{v}$ subject to 
    \[
        \begin{aligned}
            x_{u} + x_{v} &\geq 1 \text{  for all } uv \in E(G)\\
            x_{v} &\in \left[0, 1\right]  \text{  for all } v \in V(G)
        \end{aligned}
    \]
\end{proposition}

\section{Solving LP}
\subsection{Definition}
\begin{definition}[Standard form]
    We want to maximize $\sum_{j = 1}^{n} c_{j}x_{j}$ subject to 
    \[
        \begin{aligned}
            \sum_{j = 1}^{n}a_{ij}x_{j} &\leq b_{i}, \ \forall 1 \leq i \leq m\\
            x_{j} &\geq 0 \ \forall 1 \leq j \leq m
        \end{aligned}
    \]
    Standard form is to maximize $\transpose{C}x$ subject to $Ax \leq b$ and $x \geq 0$.\\
\end{definition}

\begin{proposition}
    A Linear Program can always be set in standard form.
\end{proposition}

\subsection{Geometric Interpretation}
The set of points $x \in \R^{n}$ at which a constraint holds with equality is a hyperplane. \\
Thus, each constraint is satisfied by a closed half-space of $\R^{n}$, and the set of feasible solution is the intersection of $m$ closed half-spaces, that is, a convex polyhedron $P$. \\
A linear program can have no optimal solution : 
\begin{itemize}
    \item if the set of feasible solution is empty
    \item if for every integer $M$, there exists a feasible point $x$ such that $c \cdot x \geq M$. In this case the set of feasible solution is unbounded. 
\end{itemize}
\begin{proposition}
    Because the objective function is linear, the set $\{x \mid \transpose{C}x = t\}$ is a hyperplan for any $t \in \R$. For each $t$, this is called a level set of the objective function. \\
    Thus, if the maximum is $z^{\star}$, it is a supporting hyperplane of the $P$. So $\{x \mid \transpose{C}x = z^{\star}\}$ contains an extreme point (or vertex) of $P$. 
\end{proposition}

\begin{definition}
    A supporting hyperplane of a convex set $S$ is a hyperplane such that $S \cap H \neq \emptyset$
\end{definition}

\subsection{Algorithm}
We compute the coordinates of all the vertices of the polytope : 
\begin{itemize}
    \item There are $n$ variables and $m$ constraints. 
    \item A vertex of $P$ is a solution that satifies at least $n$ inequality constraints with equality.
\end{itemize}
\begin{algorithm}
    \caption{Naive Algorithm for Linear Programs}
    \begin{algorithmic}
        \For{each set of $n$ linearly independent constraints}
            \State Check if the points that satifies all of them is feasible and compute its value
        \EndFor
        \State Get the maximum value
    \end{algorithmic}
\end{algorithm}
This is a finite algorithm. 
\subsection{Simplex Method}
\begin{algorithm}
    \caption{Simplex Method}
    \begin{algorithmic}
        \State Start at a vertex of feasible polytope
        \While {There is a vertex of better objective value}
            \State Choose a pivot
            \State Go to next vertex
        \EndWhile
    \end{algorithmic}
\end{algorithm}

\subsection{Equational Form}
\begin{definition}
    We want to maximize $\sum_{j = 1}^{n} c_{j}x_{j}$ subject to 
    \[
        \begin{aligned}
            \sum_{j = 1}^{n}a_{ij}x_{j} &= b_{i}, \ \forall 1 \leq i \leq m\\
            x_{j} &\geq 0 \ \forall 1 \leq j \leq m
        \end{aligned}
    \]
\end{definition}
\begin{proposition}
    All linear programs can be put in equational form
\end{proposition}
\begin{proof}
    We introduce slack variables that compute the difference between the right side and the left side of the inequalities in standard form
\end{proof}

\section{Analysis}
\subsection{Complexity}
\begin{proposition}
    The simplex is very fast in practice and routinely solves linear programs. However, it is a bizarre mathematical fact that the worst-case running time of the simplex method is exponential in the input size. \\
\end{proposition}
\begin{proof}
    The number of vertices of an hypercube is exponential in the dimension. A 'squashed' version of the hypercube can be used to force the simplex method to have exponential time by going through all the vertices. 
\end{proof}

\end{document}

